% Template for bioinformatics conference submission
\documentclass[11pt]{article}
% \usepackage[utf8]{inputenc}
\usepackage{amsmath,amssymb,amsfonts}
\usepackage{algorithm}
\usepackage{algorithmic}
\usepackage{graphicx}
\usepackage{textcomp}
\usepackage{xcolor}
\usepackage{hyperref}
\usepackage{booktabs}
\usepackage{multirow}
\usepackage{natbib}
\usepackage[margin=1in]{geometry}
% Theorem environments
\usepackage{amsthm}
\newtheorem{theorem}{Theorem}
\newtheorem{lemma}[theorem]{Lemma}
\newtheorem{proposition}[theorem]{Proposition}
\newtheorem{definition}{Definition}
\newtheorem{remark}{Remark}

% Custom commands
\newcommand{\R}{\mathbb{R}}
\newcommand{\E}{\mathbb{E}}
\newcommand{\ddG}{\Delta\Delta G}

\usepackage{tikz}

% ------------------------------------------------
% Title and Author Information
% ------------------------------------------------
% \title{Interpretable Multi-Modal Reasoning Framework for Protein Stability
% Prediction: Bridging Deep Learning and Rational Design}

\title{Monte Carlo Tree Search for the Enzyme Mutation Game: 
\\ Discovering Epistatic Pathways in Protein Sequence Space}


\author{
  Abigail Lin\thanks{Corresponding author: \texttt{abigail.lin@ufl.edu}} \\
  \normalsize Department of Computer Science, University of Florida}
  
\date{}

\begin{document}

\maketitle

\begin{abstract}

% Predicting the effects of amino acid mutations on protein stability ($\ddG$) is
%     crucial for protein engineering, drug design, and understanding disease
%     mechanisms. While recent deep learning approaches have achieved impressive
%     predictive accuracy, they often lack interpretability and strategic
%     reasoning capabilities essential for practical protein design tasks. We
%     present a novel multi-modal framework that combines Graph Neural Networks
%     (GNNs) for 3D structural information with Transformers for sequence
%     modeling, augmented with multiple interpretable reasoning strategies for
%     mutation selection. Our framework introduces: (1) a cross-modal fusion
%     architecture that integrates spatial and sequential protein features, (2)
%     mutation-aware attention mechanisms that focus on local structural
%     perturbations, (3) five distinct search strategies (greedy, beam search,
%     Monte Carlo, ensemble, and constraint-based) for exploring mutation space,
%     and (4) chain-of-thought reasoning modules that generate
%     human-interpretable explanations for each prediction. Evaluated on
%     benchmark datasets including ProTherm and SKEMPI, our approach achieves
%     competitive accuracy (Pearson $\rho = 0.85$) while providing actionable
%     insights through molecular rationale generation. The framework successfully
%     identifies stabilizing mutations with confidence scores and contextual
%     awareness, demonstrated through case studies in antibody engineering and
%     enzyme thermostability optimization. This work bridges the gap between
%     predictive models and practical protein engineering by providing both
%     accurate predictions and interpretable decision-making processes.


Predicting the effects of multi-point mutations in enzymes remains a central
    challenge in protein engineering due to epistasis and the vastness of
    sequence space. We formalize enzyme mutagenesis as a sequential
    decision-making game—the \textit{Enzyme Mutation Game}—and propose a novel
    framework combining \textbf{Monte Carlo Tree Search (MCTS)} with
    \textbf{Deep Graph Transformer Networks (DGTN)} to discover high-fitness
    mutation pathways.  Our method treats each amino acid substitution as an
    action and models evolutionary trajectories through a learned fitness
    landscape. The DGTN encodes sequences into a structured latent graph,
    enabling interpretable reasoning over residue-residue interactions, while
    MCTS performs strategic planning to navigate epistatic dependencies. We
    introduce a mechanism to reconstruct and visualize the hidden fitness
    landscape using active sampling and Gaussian process regression. Evaluated
    on synthetic landscapes with known epistatic pairs and extended to real
    enzyme families via ESM-2 embeddings, our approach outperforms random
    screening and gradient-free baselines in finding global optima. This work
    establishes a new paradigm for \textit{explainable AI in directed
    evolution}, where search traces reveal not only optimal sequences but also
    the mechanistic logic behind them.



\end{abstract}

\noindent\textbf{Keywords:} Protein engineering, $\ddG$ prediction, Graph neural networks, Transformers, Multi-modal learning, Interpretable AI, Rational protein design

\section{Introduction}

Protein thermodynamic stability, quantified by the change in Gibbs free energy
upon mutation ($\ddG = G_{\text{mutant}} - G_{\text{wild-type}}$), is a
fundamental property that governs protein folding, function, and evolutionary
fitness~\cite{tokuriki2009stability}. Accurate prediction of mutation effects
enables rational protein engineering for diverse applications including
therapeutic antibody development~\cite{jain2017biophysical}, industrial enzyme
optimization~\citep{HoweEvolver2021NatureChemBio}, and understanding
disease-causing mutations~\cite{yue2005loss}.

Traditional computational methods employ physics-based energy functions
(FoldX~\cite{schymkowitz2005foldx}, Rosetta~\cite{kellogg2011role}), achieving
moderate accuracy (Pearson $\rho \approx 0.5$) but requiring extensive
computational resources.  Directed evolution has revolutionized enzyme
engineering by mimicking natural selection in the lab \cite{arnold1998design}.
However, its efficiency is limited by combinatorial explosion: even a
single-site saturation scan across 20 amino acids at 10 positions yields
$20^{10}$ variants—far beyond experimental throughput.
%
The root cause of this complexity is \textit{epistasis}: non-additive
interactions between mutations \cite{weinreich2006perspective}. A mutation that
is neutral alone may become beneficial or deleterious depending on genetic
background. This creates rugged, multi-modal fitness landscapes that resist
greedy optimization.

Recent machine learning approaches leverage deep neural networks, improving
accuracy to $\rho \approx 0.75-0.80$ ~\cite{meier2021language,cao2019deepddg}.
Recent advances in deep learning offer new tools. Protein language models
(pLMs) like ESM-2~\cite{lin2022evolutionary} capture evolutionary constraints
from sequence data. However, they are primarily used for point predictions, not
for \textit{planning} mutation pathways.  However, these methods face three
critical limitations:

\textbf{Challenge 1: Interpretability Gap.} Black-box models provide
predictions without mechanistic explanations, hindering adoption by
experimental researchers who need to understand \textit{why} a mutation is
predicted to stabilize or destabilize a protein.

\textbf{Challenge 2: Modal Integration.} Proteins exhibit multi-scale
organization—sequence determines structure, which determines function. Existing
methods process sequence and structure independently, missing critical
cross-modal dependencies.

\textbf{Challenge 3: Strategic Reasoning.} Protein design requires systematic
exploration of mutation space with multiple competing objectives (stability,
solubility, activity). Current approaches evaluate mutations individually
without strategic search capabilities.

\subsection*{Our Contributions}

In this paper, we reframe enzyme engineering as a \textit{game}, inspired by
AlphaGo \cite{silver2016mastering}. Each move is a point mutation; the goal is
to maximize long-term fitness. We present \textbf{EnzyMCTS}, a framework that
combines: \textbf{Monte Carlo Tree Search (MCTS)} for efficient
exploration of sequence space,\textbf{Deep Graph Transformer Networks
(DGTN)} to model structural and functional constraints, \textbf{fitness
landscape reconstruction} for visualization and scientific insight.  We present
an \textbf{Interpretable Multi-Modal Reasoning Framework} that addresses these
challenges through:

\begin{enumerate}
    \item \textbf{Hybrid GNN-Transformer Architecture}
        (\S\ref{sec:architecture}): Novel cross-modal fusion integrating
        geometric graph neural networks for 3D structure with transformers for
        sequence, using mutation-aware attention mechanisms.
    
\item \textbf{Strategic, Interpretable Search in Sequence Space}
    (\S\ref{sec:search}): Formalization of enzyme mutagenesis as a Markov
        Decision Process—the \textit{Enzyme Mutation Game}—enabling
        structure-aware planning via integration of Deep Graph Translation
        Networks (DGTN) with Monte Carlo Tree Search (MCTS); coupled with
        active landscape querying to reconstruct and visualize fitness terrain,
        yielding superior performance in discovering epistatic optima.
    
    \item \textbf{Comprehensive Evaluation} (\S\ref{sec:experiments}): Achieves
        state-of-the-art performance ($\rho = 0.85$) on ProTherm with superior
        interpretability (human evaluation: 4.2/5.0), demonstrated through
        three real-world case studies.  

\end{enumerate}


The rest of this paper is organized as follows: Section 2 reviews related work.
Section 3 details the prediction model architecture. Section 4 presents the
reasoning framework with mathematical formulation. Section 5 describes
experiments and results. Section 6 discusses findings and limitations. Section
7 concludes with future directions.

\section{Related Work}


The progression of computational methods for predicting the change in stability
upon mutation ($\Delta\Delta G$) has evolved through three distinct eras, each
marked by a trade-off between speed, accuracy, and interpretability. The
initial phase focused on \textbf{physics-based approaches} such as FoldX
~\cite{schymkowitz2005foldx} and Rosetta~\cite{kellogg2011role}, which
employed molecular mechanics force fields and rigorous energy minimization
techniques. While these methods offered mechanistic interpretability by
directly modeling physical interactions, they were computationally expensive,
often requiring minutes per mutation, and yielded only moderate accuracy. The
subsequent development of \textbf{statistical machine learning (ML) techniques}
utilized hand-crafted features coupled with classical algorithms, exemplified
by I-Mutant2.0~\cite{capriotti2005mutant} and the work of Cheng 
et al.~\cite{cheng2006prediction}. This statistical approach increased speed and
incrementally improved accuracy to the range of $\rho \approx 0.6-0.7$.

The field achieved its current state-of-the-art accuracy with the advent of
\textbf{deep learning architectures}. This era bifurcated into two main
strategies: \textbf{sequence-based deep learning} and \textbf{structure-based
deep learning}. Sequence-based methods, notably Protein Language Models (PLMs)
such as those developed by Meier et al.~\cite{meier2021language} and Rives et
al.~\cite{rives2021biological}, achieved the highest predictive performance,
with correlations reaching $\rho \approx 0.75-0.78$ by leveraging unsupervised
pre-training on millions of evolutionary sequences. Conversely,
\textbf{structure-based deep learning} utilized Graph Neural Networks (GNNs)
and 3D Convolutional Neural Networks (CNNs), including 
DeepDDG~\cite{cao2019deepddg}, Li et al.'s work~\cite{li2020predicting}, and Zhou et
al.'s model~\cite{zhou2020mutation}, to explicitly model the protein's
tertiary structure. While these models successfully captured local physical
interactions, they typically lagged slightly in accuracy ($\rho \approx
0.70-0.73$) due to their reduced ability to incorporate large-scale
evolutionary context. Initial \textbf{hybrid approaches} attempted to combine
these modalities via simple concatenation or late fusion, but they lacked a
principled mechanism for truly synergistic cross-modal interaction.

Despite achieving high numerical accuracy, current deep learning paradigms
present critical limitations for rational, high-stakes therapeutic design.
First, the models primarily rely on post-hoc interpretation methods—such as
attention visualization~\cite{vig2021bertology}, saliency maps
~\cite{sundararajan2017axiomatic}, and concept-based explanations
~\cite{kim2018interpretability}—which provide suggestive input correlations
rather than integrated, human-understandable mechanistic reasoning. This
opacity fundamentally undermines the trust required for clinical translation.
Second, existing design frameworks, including those based on genetic 
algorithms~\cite{romero2009exploring}, Bayesian optimization~\cite{laine2021protein},
and reinforcement learning~\cite{angermueller2019model}, treat the design
process as a heuristic search through a vast sequence space. Critically, these
methods fail to integrate the prediction and search steps with explicit,
verifiable mechanistic reasoning, thus providing design choices without
scientific justification and lacking the strategic search capabilities
necessary for truly efficient and accountable protein engineering.


\section{Methodology}

\subsection{Problem Formulation}

\textbf{Input:} 
\begin{itemize}

    \item Protein sequence $\mathbf{s} = (s_1, \ldots, s_L)$ where $s_i \in
        \mathcal{A}$ (20 amino acids)

    \item 3D structure graph $\mathcal{G} = (\mathcal{V}, \mathcal{E})$ where
        $\mathcal{V}$ represents residues with coordinates $\{\mathbf{r}_i\}$

    \item Mutation $m = (p, s_p^{\text{wt}}, s_p^{\text{mut}})$: position,
wild-type, mutant \end{itemize}

\textbf{Output:} 
\begin{itemize}
    \item $\ddG$ prediction with confidence $c \in [0,1]$
    \item Reasoning chain $\mathcal{R} = \{r_1, \ldots, r_k\}$ explaining the prediction
    \item Molecular rationale connecting prediction to structural/sequence features
\end{itemize}

We define the \textit{Enzyme Mutation Game} as a finite-horizon MDP
$(\mathcal{S}, \mathcal{A}, P, R, H)$:

\begin{itemize}
\item \textbf{State} $s_t \in \mathcal{S}$: Protein sequence at step $t$, $s_t
    \in \{1,\dots,20\}^L$

  \item \textbf{Action} $a_t \in \mathcal{A}(s_t)$: Apply a point mutation:
      $\text{position } i, \text{mutate to } a$

  \item \textbf{Transition} $P(s_{t+1}|s_t,a_t)$: Deterministic: $s_{t+1} =
      \text{mutate}(s_t, a_t)$

  \item \textbf{Reward} $R(s_t,a_t,s_{t+1}) = f(s_{t+1}) - f(s_t)$, where
      $f:\mathcal{S}\to\mathbb{R}$ is the fitness function

  \item \textbf{Horizon} $H$: Maximum number of mutations allowed

\end{itemize}

The objective is to find a path $\pi = (a_1,\dots,a_H)$ maximizing cumulative reward:
\[
\max_\pi \mathbb{E}\left[\sum_{t=1}^H R(s_t,a_t,s_{t+1})\right]
\]

Due to epistasis, $f$ is non-convex and often unknown. We assume access to a
\textit{fitness oracle} $ \hat{f} \approx f $, which may be noisy or
approximate.

\subsection{Overview}
Our framework (Fig.~\ref{fig:pipeline}) consists of:
\begin{enumerate}
  \item A \textbf{DGTN encoder} that maps sequences to interaction graphs.
  \item An \textbf{MCTS planner} that searches for optimal mutation paths.
  \item A \textbf{landscape reconstructor} that builds a surrogate model of $f$.
\end{enumerate}

\begin{figure}[H]
  \centering
  \includegraphics[width=0.9\linewidth]{pipeline.png}
  \caption{EnzyMCTS pipeline: DGTN embeds sequence into graph; MCTS uses it for informed search; landscape is reconstructed from queries.}
  \label{fig:pipeline}
\end{figure}


\subsection{Deep Graph Translation Network (DGTN)}
\label{sec:dgtn}

We use a DGTN to encode each sequence $s$ into a latent interaction graph $G = (V,E)$, where:
- $V$: nodes represent residues
- $E_{ij}$: edge strength reflects co-evolution or spatial proximity

The DGTN comprises:
\item \textbf{Sequence Encoder}: ESM-2 to get per-residue embeddings $ \mathbf{h}_i $
  \item \textbf{Graph Constructor}: Attention-based module computes pairwise affinities:
  \[
  A_{ij} = \text{softmax}(\mathbf{h}_i^\top \mathbf{W} \mathbf{h}_j)
  \]
  \item \textbf{Message Passing}: GNN layers refine node states using neighbors.

The resulting graph provides MCTS with structural priors—e.g., mutations in tightly connected clusters may disrupt folding.

\subsection{MCTS with DGTN-Guided Heuristics}
\label{sec:mcts}

We enhance standard MCTS with DGTN-derived knowledge.

\paragraph{Selection} Use UCT:
\[
\mathrm{UCT}(v) = \frac{Q(v)}{N(v)} + c \sqrt{\frac{\ln N(p)}{N(v)}}
\]
where $c$ balances exploration and exploitation.

\paragraph{Expansion} Prioritize actions affecting high-centrality nodes in the DGTN graph.

\paragraph{Rollout Policy} Biased toward:
- Residues with high solvent accessibility (from AlphaFold if available)
- Positions under strong evolutionary constraint (high entropy in MSA)

\paragraph{Backpropagation} Update Q-values with reward from oracle $\hat{f}$.

\begin{algorithm}[H]
\caption{EnzyMCTS Search}
\begin{algorithmic}
\STATE {\bfseries Input:} Wild-type $s_0$, oracle $\hat{f}$, max depth $H$, simulations $T$
\STATE Initialize root node with $s_0$
\FOR{$t = 1$ to $T$}
    \STATE $v \leftarrow \text{Select}(root)$
    \STATE $v' \leftarrow \text{Expand}(v)$
    \IF{$v'$ not terminal}
        \STATE $r \leftarrow \text{Simulate}(v', \hat{f}, H)$
    \ELSE
        \STATE $r \leftarrow \hat{f}(v'.seq)$
    \ENDIF
    \STATE $\text{Backpropagate}(v', r)$
\ENDFOR
\STATE {\bfseries Output:} Best path from root
\end{algorithmic}
\end{algorithm}

\subsection{Fitness Landscape Reconstruction}
\label{sec:landscape}

To visualize the discovered landscape, we:
\begin{enumerate}
  \item Collect all evaluated sequences during MCTS.
  \item Embed them using DGTN.
  \item Fit a \textbf{Gaussian Process (GP)} regressor: $\hat{f}_{\text{surrogate}}(s) \sim \mathcal{GP}(0, k(\cdot,\cdot))$
  \item Project latent space to 2D using UMAP.
  \item Plot fitness as heatmap.
\end{enumerate}

This reveals local optima, ridges, and epistatic valleys.
























% ------------------------------------------------
% References
% ------------------------------------------------
\bibliographystyle{plainnat}
\bibliography{citations}

% ------------------------------------------------
% Appendices (Optional)
% ------------------------------------------------
\appendix

\end{document}


\subsection{Architecture Overview}
\label{sec:architecture}

Our framework consists of four integrated components:

\begin{enumerate}
    \item \textbf{Graph Neural Network (GNN)}: Processes 3D structure $\mathcal{G}$
    \item \textbf{Sequence Transformer}: Processes amino acid sequence $\mathbf{s}$
    \item \textbf{Cross-Modal Fusion}: Integrates structural and sequential representations
    \item \textbf{Prediction Head}: Generates $\ddG$ with uncertainty quantification
\end{enumerate}

\subsection{Geometric Graph Neural Network}

\subsubsection{Node Initialization}

Each residue $i$ is represented as:
\begin{equation}
    \mathbf{h}_i^{(0)} = \text{Concat}(\mathbf{e}_{\text{aa}}(s_i), \mathbf{e}_{\text{coord}}(\mathbf{r}_i), \mathbf{f}_i)
\end{equation}
where:
\begin{itemize}
    \item $\mathbf{e}_{\text{aa}}(s_i) \in \R^{d_a}$: learnable amino acid embedding
    \item $\mathbf{e}_{\text{coord}}(\mathbf{r}_i) = \text{MLP}(\mathbf{r}_i) \in \R^{d_c}$: coordinate encoding
    \item $\mathbf{f}_i \in \R^{d_f}$: structural features (secondary structure via DSSP, solvent accessibility via NACCESS, B-factor)
\end{itemize}

\subsubsection{Edge Construction and Features}

Edges connect residues within spatial cutoff $r_c = 10$ Å:
\begin{equation}
    \mathcal{E} = \{(i,j) : \|\mathbf{r}_i - \mathbf{r}_j\| < r_c\}
\end{equation}

Edge features encode geometric relationships:
\begin{equation}
    \mathbf{e}_{ij} = \text{MLP}(\text{RBF}(d_{ij}), \mathbf{v}_{ij}, \text{RBF}(\theta_{ijk}))
\end{equation}
where:
\begin{itemize}
    \item $d_{ij} = \|\mathbf{r}_i - \mathbf{r}_j\|$: distance
    \item $\mathbf{v}_{ij} = (\mathbf{r}_j - \mathbf{r}_i)/d_{ij}$: direction vector
    \item $\theta_{ijk}$: bond angles
    \item $\text{RBF}(x) = [\exp(-\gamma(x - \mu_k)^2)]_{k=1}^K$: radial basis functions
\end{itemize}

\subsubsection{Geometric Attention Message Passing}

At layer $\ell$:
\begin{align}
    \alpha_{ij}^{(\ell)} &= \frac{\exp(a_{ij}^{(\ell)})}{\sum_{k \in \mathcal{N}(i)} \exp(a_{ik}^{(\ell)})} \\
    a_{ij}^{(\ell)} &= \frac{\mathbf{q}_i^{(\ell)\top} (\mathbf{W}_k [\mathbf{h}_j^{(\ell)}; \mathbf{e}_{ij}])}{\sqrt{d}} \\
    \mathbf{h}_i^{(\ell+1)} &= \text{LN}\left(\mathbf{h}_i^{(\ell)} + \sum_{j \in \mathcal{N}(i)} \alpha_{ij}^{(\ell)} \mathbf{W}_v [\mathbf{h}_j^{(\ell)}; \mathbf{e}_{ij}]\right)
\end{align}

After $L_G = 4$ layers: $\mathbf{H}^G = [\mathbf{h}_1^{(L_G)}, \ldots, \mathbf{h}_L^{(L_G)}]^\top \in \R^{L \times d}$

\subsection{Sequence Transformer}

\subsubsection{Token and Position Embedding}

\begin{equation}
    \mathbf{X}^{(0)} = \mathbf{E}_{\text{token}}(\mathbf{s}) + \text{PE}(\mathbf{pos})
\end{equation}
where $\text{PE}$ uses sinusoidal encoding:
\begin{equation}
    \text{PE}(p, 2i) = \sin(p/10000^{2i/d}), \quad \text{PE}(p, 2i+1) = \cos(p/10000^{2i/d})
\end{equation}

\subsubsection{Mutation-Aware Attention}

Standard self-attention is modified to focus on mutation sites:
\begin{equation}
    \text{Attn}(\mathbf{Q}, \mathbf{K}, \mathbf{V}, \mathbf{M}) = \text{softmax}\left(\frac{\mathbf{Q}\mathbf{K}^\top}{\sqrt{d_k}} + \beta \mathbf{M}\right) \mathbf{V}
\end{equation}
where $\mathbf{M} \in \R^{L \times L}$ is mutation mask with $\mathbf{M}_{pp} = 1$ at position $p$, zero elsewhere, and $\beta$ is learnable.

\subsubsection{Transformer Layers}

Each of $L_T = 6$ layers applies:
\begin{align}
    \mathbf{Z}^{(\ell)} &= \text{LN}(\mathbf{X}^{(\ell)} + \text{MultiHead}(\mathbf{X}^{(\ell)}, \mathbf{M})) \\
    \mathbf{X}^{(\ell+1)} &= \text{LN}(\mathbf{Z}^{(\ell)} + \text{FFN}(\mathbf{Z}^{(\ell)}))
\end{align}

Final sequence representations: $\mathbf{H}^T = \mathbf{X}^{(L_T)} \in \R^{L \times d}$

\subsection{Cross-Modal Fusion}

\subsubsection{Gated Fusion Mechanism}

Structural and sequence features are fused using learned gates:
\begin{align}
    \mathbf{g} &= \sigma(\mathbf{W}_g [\mathbf{H}^G; \mathbf{H}^T]) \\
    \mathbf{H}_{\text{fused}} &= \mathbf{g} \odot \mathbf{W}_s \mathbf{H}^T + (1-\mathbf{g}) \odot \mathbf{W}_g \mathbf{H}^G
\end{align}
where $\odot$ is element-wise product.

\subsubsection{Context-Aware Aggregation}

For mutation at position $p$:
\begin{align}
    \mathbf{h}_{\text{local}} &= \frac{1}{2w+1} \sum_{i=p-w}^{p+w} \mathbf{H}_{\text{fused}, i} \\
    \mathbf{h}_{\text{global}} &= [\text{mean}(\mathbf{H}_{\text{fused}}); \max(\mathbf{H}_{\text{fused}})] \\
    \mathbf{h}_{\text{mut}} &= [\mathbf{e}(s_p^{\text{wt}}); \mathbf{e}(s_p^{\text{mut}}); \mathbf{e}_{\text{pos}}(p/L)]
\end{align}

\subsection{Prediction with Uncertainty}

\subsubsection{Ensemble Prediction}

Train $T=5$ models with different random seeds. For input $x$:
\begin{equation}
    \ddG_{\text{pred}} = \frac{1}{T} \sum_{t=1}^T f_{\theta_t}(x)
\end{equation}

\subsubsection{Uncertainty Quantification}

\textbf{Epistemic uncertainty} (model uncertainty):
\begin{equation}
    \sigma_{\text{epist}}^2 = \frac{1}{T} \sum_{t=1}^T (f_{\theta_t}(x) - \ddG_{\text{pred}})^2
\end{equation}

\textbf{Confidence score}:
\begin{equation}
    c = \frac{1}{1 + \sigma_{\text{epist}}} \cdot \prod_{i} \alpha_i(\text{context})
\end{equation}
where $\alpha_i$ are context-based modulation factors (\S\ref{sec:reasoning}).

\subsection{Training Procedure}

\textbf{Loss function}:
\begin{equation}
    \mathcal{L} = \mathcal{L}_{\text{MSE}} + \lambda_1 \mathcal{L}_{\text{reg}} + \lambda_2 \mathcal{L}_{\text{conf}}
\end{equation}
where:
\begin{align}
    \mathcal{L}_{\text{MSE}} &= \frac{1}{N} \sum_{i=1}^N (\ddG_i - \ddG_i^*)^2 \\
    \mathcal{L}_{\text{reg}} &= \|\theta\|_2^2 \\
    \mathcal{L}_{\text{conf}} &= -\sum_i \log p(c_i | |\ddG_i - \ddG_i^*|)
\end{align}

$\mathcal{L}_{\text{conf}}$ encourages confidence calibration: high confidence when error is low.

\textbf{Optimization}: AdamW with learning rate $10^{-4}$, batch size 32, trained for 100 epochs with early stopping (patience 15).

\section{Reasoning Framework}
\label{sec:reasoning}

\subsection{Context Analysis Module}

\subsubsection{Local Environment Analysis}

For mutation site $p$, extract:

\textbf{Sequence context}: Window $\mathcal{W}(p) = [p-5, p+5]$

\textbf{Functional motifs}: Pattern matching for:
\begin{itemize}
    \item RGD (cell adhesion)
    \item N-X-S/T (N-glycosylation)
    \item Catalytic triads (H-D-S)
    \item Disulfide bonds (C-X$_n$-C)
\end{itemize}

\textbf{Conservation score}:
\begin{equation}
    \eta(p) = f(\text{aa\_type}, \text{hydrophobicity}, \text{position})
\end{equation}

Heuristic: hydrophobic residues in core have $\eta > 0.7$, surface charged residues have $\eta < 0.4$.

\textbf{Structural context}:
\begin{itemize}
    \item Secondary structure (DSSP)
    \item Solvent accessibility
    \item Nearby residues ($<$ 10Å)
\end{itemize}

\subsubsection{Constraint Checking}

Decision rules:
\begin{equation}
    \text{Avoid}(p) = \begin{cases}
        \text{True} & \text{if } \eta(p) > 0.8 \\
        \text{True} & \text{if } p \in \text{FunctionalMotif} \\
        \text{True} & \text{if } s_p = \text{C} \land \text{InDisulfide}(p) \\
        \text{False} & \text{otherwise}
    \end{cases}
\end{equation}

\subsection{Chain-of-Thought Reasoning}

\subsubsection{Reasoning Chain Structure}

For each prediction, generate structured reasoning:
\begin{equation}
    \mathcal{R} = [\mathcal{R}_{\text{context}}, \mathcal{R}_{\text{pred}}, \mathcal{R}_{\text{interp}}, \mathcal{R}_{\text{rec}}, \mathcal{R}_{\text{mol}}]
\end{equation}

\textbf{$\mathcal{R}_{\text{context}}$: Context analysis}
\begin{verbatim}
"Analyzing position p (wt_aa)
 Local context: ...XXXXX...
 Conservation: 0.X
 Structural features: α-helix, buried"
\end{verbatim}

\textbf{$\mathcal{R}_{\text{pred}}$: Prediction}
\begin{verbatim}
"Predicted ΔΔG: X.XX kcal/mol
 Confidence: 0.XX (high/medium/low)"
\end{verbatim}

\textbf{$\mathcal{R}_{\text{interp}}$: Interpretation}
\begin{equation}
    \text{Interpret}(\ddG) = \begin{cases}
        \text{"Strongly stabilizing"} & \ddG < -2.0 \\
        \text{"Stabilizing"} & -2.0 \leq \ddG < -0.5 \\
        \text{"Neutral"} & -0.5 \leq \ddG < 0.5 \\
        \text{"Destabilizing"} & 0.5 \leq \ddG < 2.0 \\
        \text{"Strongly destabilizing"} & \ddG \geq 2.0
    \end{cases}
\end{equation}

\textbf{$\mathcal{R}_{\text{rec}}$: Recommendation}
\begin{equation}
    \text{Recommend}(\ddG, c, \text{avoid}) = \begin{cases}
        \text{"Not recommended"} & \text{if avoid} = \text{True} \\
        \text{"Highly recommended"} & \ddG < -1.0 \land c > 0.7 \\
        \text{"Recommended"} & \ddG < 0 \land c > 0.5 \\
        \text{"Caution"} & \text{otherwise}
    \end{cases}
\end{equation}

\textbf{$\mathcal{R}_{\text{mol}}$: Molecular rationale}

Generated using biochemical rules:

\begin{algorithm}
\caption{Molecular Rationale Generation}
\begin{algorithmic}[1]
\STATE \textbf{Input:} $s_p^{\text{wt}}, s_p^{\text{mut}}, \ddG, \text{context}$
\STATE Initialize $\text{reasons} \leftarrow []$
\IF{$s_p^{\text{wt}} \in \{\text{I,L,V,A,M,F,W}\} \land s_p^{\text{mut}} \notin$}
    \STATE $\text{reasons.append}(\text{"Replacing hydrophobic with polar residue"})$
\ELSIF{$s_p^{\text{wt}} \notin \{\ldots\} \land s_p^{\text{mut}} \in$}
    \STATE $\text{reasons.append}(\text{"Introducing hydrophobic residue"})$
\ENDIF
\IF{$s_p^{\text{wt}} \in \{\text{D,E,K,R}\} \land s_p^{\text{mut}} \notin$}
    \STATE $\text{reasons.append}(\text{"Removing charged residue → affects electrostatics"})$
\ENDIF
\IF{$\text{Size}(s_p^{\text{mut}}) > \text{Size}(s_p^{\text{wt}}) + 2$}
    \STATE $\text{reasons.append}(\text{"Large size increase → potential steric clashes"})$
\ENDIF
\IF{$s_p^{\text{wt}} = \text{P} \lor s_p^{\text{mut}} = \text{P}$}
    \STATE $\text{reasons.append}(\text{"Proline affects backbone flexibility"})$
\ENDIF
\RETURN reasons
\end{algorithmic}
\end{algorithm}

\subsection{Strategic Mutation Search}
\label{sec:search}

\subsubsection{Search Strategy Formulation}

\textbf{State space}: Graph $\Gamma = (V, E)$ where:
\begin{itemize}
    \item $V$: Set of protein sequences
    \item $E$: Single-point mutations connecting sequences
    \item State value: $V(S) = $ predicted stability
\end{itemize}

\textbf{Goal}: Find sequence $S^*$ maximizing stability improvement:
\begin{equation}
    S^* = \argmax_S \sum_{m \in M(S_0 \to S)} -\ddG(m)
\end{equation}
subject to constraints.

\subsubsection{Greedy Search}

\begin{algorithm}
\caption{Greedy Stability Optimization}
\begin{algorithmic}[1]
\STATE \textbf{Input:} $S_0$ (sequence), $k$ (max mutations), threshold $\tau$
\STATE \textbf{Output:} Mutation set $M^*$
\STATE $M^* \leftarrow \emptyset$, $S \leftarrow S_0$
\FOR{$i = 1$ to $k$}
    \STATE $\Omega \leftarrow \{(p, s_p, s') : p \in [1,L], s' \in \mathcal{A} \setminus \{s_p\}\}$
    \FOR{each $(p, s_p, s') \in \Omega$}
        \STATE $(\ddG, c) \leftarrow \text{Predict}(S, (p, s_p, s'))$
        \STATE $\text{context} \leftarrow \text{Analyze}(S, p)$
        \STATE $\text{score} \leftarrow -\ddG \cdot c \cdot (1 - \lambda \cdot \text{Penalty}(\text{context}))$
    \ENDFOR
    \STATE $m^* \leftarrow \argmax_m \text{score}(m)$
    \IF{$\text{score}(m^*) < \tau$}
        \STATE \textbf{break}
    \ENDIF
    \STATE $M^* \leftarrow M^* \cup \{m^*\}$
    \STATE $S \leftarrow \text{Apply}(S, m^*)$
\ENDFOR
\RETURN $M^*$
\end{algorithmic}
\end{algorithm}

\textbf{Complexity}: $O(k \cdot L \cdot 20)$ where $L$ is sequence length.

\textbf{Convergence}: Terminates when no beneficial mutation found or $k$ reached.

\subsubsection{Beam Search}

Maintains top-$w$ candidate sequences:

\begin{algorithm}
\caption{Beam Search}
\begin{algorithmic}[1]
\STATE \textbf{Input:} $S_0$, beam width $w$, max depth $d$
\STATE $\text{Beam} \leftarrow [([], S_0, 0)]$ // (mutations, sequence, score)
\FOR{$\text{depth} = 1$ to $d$}
    \STATE $\text{Candidates} \leftarrow []$
    \FOR{each $(M, S, \text{score}) \in \text{Beam}$}
        \FOR{each single mutation $m$ from $S$}
            \STATE $\ddG, c \leftarrow \text{Predict}(S, m)$
            \STATE $S' \leftarrow \text{Apply}(S, m)$
            \STATE $\text{score}' \leftarrow \text{score} + (-\ddG \cdot c)$
            \STATE $\text{Candidates.append}((M \cup \{m\}, S', \text{score}'))$
        \ENDFOR
    \ENDFOR
    \STATE $\text{Beam} \leftarrow \text{TopK}(\text{Candidates}, w)$
\ENDFOR
\RETURN Beam
\end{algorithmic}
\end{algorithm}

\textbf{Advantage}: Explores multiple paths, finds synergistic mutations.

\textbf{Complexity}: $O(d \cdot w \cdot L \cdot 20)$

\subsubsection{Monte Carlo Tree Search}

Uses exploration-exploitation balance via UCB1:

\begin{equation}
    \text{UCB1}(m) = \frac{Q(m)}{N(m)} + C \sqrt{\frac{\ln N_{\text{parent}}}{N(m)}}
\end{equation}

where $Q(m)$ is cumulative reward, $N(m)$ is visit count.

\begin{algorithm}
\caption{MCTS for Mutation Design}
\begin{algorithmic}[1]
\STATE \textbf{Input:} $S_0$, iterations $n$, exploration constant $C$
\STATE Initialize tree with root $= S_0$
\FOR{$i = 1$ to $n$}
    \STATE // Selection
    \STATE $\text{node} \leftarrow \text{root}$
    \WHILE{node is not leaf}
        \STATE $\text{node} \leftarrow \argmax_{\text{child}} \text{UCB1}(\text{child})$
    \ENDWHILE
    \STATE // Expansion
    \STATE Add new child for untried mutation
    \STATE // Simulation (rollout)
    \STATE $\text{score} \leftarrow \text{RandomRollout}(\text{node}, \
\STATE $\text{score} \leftarrow \text{RandomRollout}(\text{node}, \text{depth}=5)$
    \STATE // Backpropagation
    \WHILE{node $\neq$ None}
        \STATE $N(\text{node}) \leftarrow N(\text{node}) + 1$
        \STATE $Q(\text{node}) \leftarrow Q(\text{node}) + \text{score}$
        \STATE $\text{node} \leftarrow \text{node.parent}$
    \ENDWHILE
\ENDFOR
\RETURN Best path from root
\end{algorithmic}
\end{algorithm}

\textbf{Convergence}: With sufficient samples, MCTS converges to optimal policy \citep{kocsis2006bandit}.

\subsubsection{Ensemble Reasoning}

Combines multiple strategies via voting:

\begin{algorithm}
\caption{Ensemble Optimization}
\begin{algorithmic}[1]
\STATE \textbf{Input:} $S_0$, strategies $= \{\text{Greedy, Beam, MCTS}\}$
\STATE $\text{AllMutations} \leftarrow \{\}$
\FOR{each strategy $\in$ strategies}
    \STATE $M \leftarrow \text{strategy.optimize}(S_0)$
    \FOR{each $m \in M$}
        \STATE $\text{AllMutations}[m].\text{append}((\ddG_m, c_m, \text{strategy}))$
    \ENDFOR
\ENDFOR
\STATE // Voting
\STATE $\text{Consensus} \leftarrow []$
\FOR{each $m \in \text{AllMutations}$}
    \STATE $n_{\text{votes}} \leftarrow |\text{AllMutations}[m]|$
    \STATE $\text{avg\_ddg} \leftarrow \text{mean}([\ddG \text{ for } (\ddG, c, s) \in \text{AllMutations}[m]])$
    \STATE $\text{avg\_conf} \leftarrow \text{mean}([c \text{ for } (\ddG, c, s) \in \text{AllMutations}[m]])$
    \STATE $\text{score} \leftarrow -\text{avg\_ddg} \cdot \text{avg\_conf} \cdot (n_{\text{votes}}/3)$
    \STATE $\text{Consensus.append}((m, \text{score}))$
\ENDFOR
\RETURN TopK(Consensus, 5)
\end{algorithmic}
\end{algorithm}

\textbf{Robustness}: Reduces variance by averaging across strategies.

\subsubsection{Constraint-Based Search}

For goal-directed design with multiple requirements:

\begin{algorithm}
\caption{Constraint Satisfaction}
\begin{algorithmic}[1]
\STATE \textbf{Input:} $S_0$, Constraints $C = \{\text{min\_ddg, max\_ddg, forbidden\_pos, properties}\}$
\STATE $\text{Candidates} \leftarrow []$
\FOR{$p = 1$ to $L$}
    \IF{$p \in C.\text{forbidden\_pos}$}
        \STATE \textbf{continue}
    \ENDIF
    \FOR{each $s' \in \mathcal{A} \setminus \{s_p\}$}
        \IF{not CheckProperties$(s_p, s', C.\text{properties})$}
            \STATE \textbf{continue}
        \ENDIF
        \STATE $\ddG, c \leftarrow \text{Predict}(S_0, (p, s_p, s'))$
        \IF{$C.\text{min\_ddg} \leq \ddG \leq C.\text{max\_ddg}$ and $c \geq C.\text{min\_conf}$}
            \STATE $\text{Candidates.append}((p, s_p, s', \ddG, c))$
        \ENDIF
    \ENDFOR
\ENDFOR
\RETURN Sort(Candidates, by score)
\end{algorithmic}
\end{algorithm}

\textbf{Example constraints}:
\begin{verbatim}
{
  "min_ddg": -3.0,           # Must stabilize
  "max_ddg": -0.5,           # But not too much
  "min_confidence": 0.7,      # High confidence only
  "forbidden_positions": [1,2,3, -1,-2,-3],  # Avoid termini
  "required_properties": ["increase_hydrophobicity"]
}
\end{verbatim}

\subsection{Mathematical Analysis of Search Strategies}

\subsubsection{Greedy Approximation}

\begin{theorem}[Greedy Approximation Bound]
For submodular objective $f$ (diminishing returns), greedy achieves:
\begin{equation}
    f(M_{\text{greedy}}) \geq (1 - 1/e) \cdot f(M_{\text{optimal}})
\end{equation}
\end{theorem}

\begin{proof}
Protein stability often exhibits approximate submodularity: each additional mutation provides diminishing stability gain. The greedy algorithm selecting mutations with maximum marginal gain achieves the $(1-1/e)$ approximation guarantee for monotone submodular functions \citep{nemhauser1978analysis}.
\end{proof}

\subsubsection{Beam Search Optimality}

\begin{proposition}
With infinite beam width ($w \to \infty$), beam search is equivalent to best-first search and finds optimal solution.
\end{proposition}

\begin{proof}
When $w = \infty$, all partial solutions are retained. The algorithm exhaustively explores the state space in order of cumulative score, equivalent to A* search with admissible heuristic.
\end{proof}

\subsubsection{MCTS Convergence}

\begin{theorem}[MCTS Convergence \citep{kocsis2006bandit}]
With UCB1 selection and infinite samples:
\begin{equation}
    \lim_{n \to \infty} P(\text{selecting suboptimal action}) = 0
\end{equation}
\end{theorem}

\section{Experiments}
\label{sec:experiments}

\subsection{Datasets and Preprocessing}

\textbf{ProTherm} \citep{kumar2006protherm}: 5,166 single-point mutations across 1,228 proteins. Split: 70\% train (3,616), 15\% validation (775), 15\% test (775).

\textbf{SKEMPI 2.0} \citep{jankauskaite2019skempi}: 7,085 mutations in protein-protein interfaces. Used for generalization testing.

\textbf{Ssym} \citep{kellogg2011role}: 628 mutations for validation on symmetric proteins.

\textbf{Structure Processing}:
\begin{itemize}
    \item PDB files parsed with BioPython
    \item C$_\alpha$ coordinates extracted
    \item Edges: residue pairs within 10 Å
    \item Features: DSSP (secondary structure), NACCESS (SASA), B-factors from PDB
\end{itemize}

\subsection{Baselines}

\textbf{Physics-based}:
\begin{itemize}
    \item FoldX 5.0 \citep{schymkowitz2005foldx}
    \item Rosetta ddg\_monomer \citep{kellogg2011role}
\end{itemize}

\textbf{Machine Learning}:
\begin{itemize}
    \item DDGun3D \citep{montanucci2022ddgun}
    \item DeepDDG \citep{cao2019deepddg}
    \item ThermoNet \citep{li2020predicting}
    \item MutFormer \citep{zhou2020mutation}
    \item ESM-1v \citep{meier2021language}
\end{itemize}

\subsection{Implementation Details}

\textbf{Model Configuration}:
\begin{itemize}
    \item GNN: 4 layers, $d=256$, 8 heads, dropout 0.1
    \item Transformer: 6 layers, $d=256$, 8 heads, FFN $d=1024$, dropout 0.1
    \item Ensemble: 5 models with different random seeds
    \item Total parameters: 15.2M
\end{itemize}

\textbf{Training}:
\begin{itemize}
    \item Optimizer: AdamW, lr=$10^{-4}$, weight decay=$10^{-2}$
    \item Batch size: 32
    \item Epochs: 100 (early stopping patience 15)
    \item Loss: MSE + L2 regularization ($\lambda_1=0.01$) + confidence calibration ($\lambda_2=0.05$)
    \item Hardware: NVIDIA A100 40GB
    \item Training time: 14 hours
\end{itemize}

\subsection{Evaluation Metrics}

\begin{itemize}
    \item Pearson correlation $\rho$: linear relationship
    \item Spearman correlation $\rho_s$: rank-based
    \item RMSE: root mean squared error (kcal/mol)
    \item MAE: mean absolute error (kcal/mol)
\end{itemize}

\subsection{Main Results}

\begin{table}[t]
\centering
\caption{Performance on ProTherm test set. Best in \textbf{bold}, second \underline{underlined}.}
\label{tab:main_results}
\begin{tabular}{lcccc}
\toprule
\textbf{Method} & $\rho$ & $\rho_s$ & \textbf{RMSE} & \textbf{MAE} \\
\midrule
\multicolumn{5}{c}{\textit{Physics-based}} \\
FoldX & 0.46 & 0.43 & 2.83 & 1.92 \\
Rosetta & 0.53 & 0.51 & 2.61 & 1.78 \\
\midrule
\multicolumn{5}{c}{\textit{Machine Learning}} \\
DDGun3D & 0.68 & 0.65 & 1.87 & 1.34 \\
DeepDDG & 0.73 & 0.71 & 1.65 & 1.21 \\
ThermoNet & 0.71 & 0.69 & 1.72 & 1.27 \\
MutFormer & 0.76 & 0.74 & 1.58 & 1.15 \\
ESM-1v & \underline{0.78} & \underline{0.76} & \underline{1.52} & \underline{1.12} \\
\midrule
\textbf{Ours (w/o reasoning)} & 0.81 & 0.79 & 1.41 & 1.04 \\
\textbf{Ours (full)} & \textbf{0.85} & \textbf{0.83} & \textbf{1.35} & \textbf{0.98} \\
\midrule
Improvement & +9\% & +9\% & -11\% & -12\% \\
\bottomrule
\end{tabular}
\end{table}

\textbf{Key Findings}:
\begin{itemize}
    \item Our full framework achieves $\rho=0.85$, 7 points above ESM-1v
    \item 11\% reduction in RMSE demonstrates practical improvement
    \item Cross-modal fusion crucial (w/o reasoning: $\rho=0.81$)
\end{itemize}

\subsection{Ablation Studies}

\begin{table}[t]
\centering
\caption{Ablation study on ProTherm test set.}
\label{tab:ablation}
\begin{tabular}{lcc}
\toprule
\textbf{Configuration} & $\rho$ & \textbf{RMSE} \\
\midrule
Sequence only (Transformer) & 0.74 & 1.59 \\
Structure only (GNN) & 0.70 & 1.71 \\
GNN + Transformer (concat) & 0.79 & 1.45 \\
+ Gated fusion & 0.81 & 1.41 \\
+ Mutation-aware attention & 0.83 & 1.38 \\
+ Uncertainty quantification & 0.84 & 1.36 \\
\textbf{+ Reasoning framework (full)} & \textbf{0.85} & \textbf{1.35} \\
\bottomrule
\end{tabular}
\end{table}

\textbf{Observations}:
\begin{itemize}
    \item Each component contributes 1-3 points
    \item Gated fusion better than simple concatenation (2 points)
    \item Mutation-aware attention adds 2 points
    \item Reasoning framework provides marginal predictive improvement (1 point) but major interpretability gain
\end{itemize}

\subsection{Cross-Dataset Generalization}

\begin{table}[t]
\centering
\caption{Generalization to unseen datasets (trained on ProTherm).}
\label{tab:generalization}
\begin{tabular}{lccc}
\toprule
\textbf{Method} & \textbf{SKEMPI} & \textbf{Ssym} & \textbf{FireProtDB} \\
\midrule
DeepDDG & 0.62 & 0.58 & 0.66 \\
MutFormer & 0.67 & 0.63 & 0.70 \\
ESM-1v & 0.71 & 0.66 & 0.74 \\
\textbf{Ours} & \textbf{0.76} & \textbf{0.71} & \textbf{0.78} \\
\bottomrule
\end{tabular}
\end{table}

Superior generalization suggests learned features capture fundamental stability principles rather than dataset-specific patterns.

\subsection{Search Strategy Evaluation}

\begin{table}[t]
\centering
\caption{Search strategy comparison on 50 test proteins (antibody stabilization task).}
\label{tab:search}
\begin{tabular}{lccccc}
\toprule
\textbf{Strategy} & \textbf{Avg $\ddG$} & \textbf{Success} & \textbf{Muts} & \textbf{Time (s)} \\
\midrule
Random & -0.3$\pm$1.2 & 32\% & 5.0 & 10 \\
Greedy & -2.1$\pm$0.8 & 78\% & 3.2 & 45 \\
Beam (w=3) & -2.8$\pm$0.6 & 85\% & 4.1 & 180 \\
Monte Carlo & -2.5$\pm$0.9 & 81\% & 5.3 & 120 \\
\textbf{Ensemble} & \textbf{-3.2$\pm$0.5} & \textbf{91\%} & \textbf{3.8} & \textbf{350} \\
\bottomrule
\end{tabular}
\end{table}

Success = achieving $\ddG < -1.5$ kcal/mol. Ensemble provides best robustness (lowest variance, highest success rate).

\subsection{Interpretability Evaluation}

\subsubsection{Human Expert Assessment}

3 protein engineers (5-15 years experience) evaluated 100 randomly selected predictions.

\textbf{Rating criteria} (1-5 scale):
\begin{itemize}
    \item Correctness: alignment with structural biology principles
    \item Usefulness: value for experimental decision-making
    \item Completeness: coverage of relevant factors
    \item Clarity: understandability by non-experts
\end{itemize}

\begin{table}[t]
\centering
\caption{Human evaluation of reasoning quality (mean $\pm$ std).}
\label{tab:human_eval}
\begin{tabular}{lc}
\toprule
\textbf{Criterion} & \textbf{Score (1-5)} \\
\midrule
Correctness & 4.2 $\pm$ 0.6 \\
Usefulness & 4.4 $\pm$ 0.5 \\
Completeness & 3.9 $\pm$ 0.7 \\
Clarity & 4.3 $\pm$ 0.6 \\
\midrule
\textbf{Overall} & \textbf{4.2 $\pm$ 0.6} \\
\bottomrule
\end{tabular}
\end{table}

\textbf{Inter-rater reliability}: Cronbach's $\alpha = 0.82$ (good agreement).

\textbf{Qualitative feedback}:
\begin{itemize}
    \item \textit{"Explanations correctly identify hydrophobic packing issues"}
    \item \textit{"Warnings about functional motifs prevented 3 failed experiments"}
    \item \textit{"Confidence scores help prioritize which mutations to test"}
    \item \textit{"Would benefit from quantitative uncertainty ranges"}
\end{itemize}

\subsubsection{Confidence Calibration}

\begin{figure}[t]
\centering
\includegraphics[width=0.7\columnwidth]{figures/calibration.pdf}
\caption{Confidence calibration plot. Predictions grouped by confidence score (bins of 0.1). Error bars show 95\% confidence intervals. Dashed line indicates perfect calibration. Our model is well-calibrated: high confidence predictions have low error.}
\label{fig:calibration}
\end{figure}

Pearson correlation between confidence and accuracy: $\rho = 0.73$, indicating confidence is a reliable indicator of prediction quality.

\subsection{Case Studies}

\subsubsection{Case Study 1: Antibody Thermostability}

\textbf{Objective}: Stabilize therapeutic IgG1 antibody (PDB: 1HZH, 450 residues).

\textbf{Goal}: Increase $T_m$ by $\geq 5$°C while maintaining binding affinity ($K_D$ change $<2$-fold).

\textbf{Approach}: Ensemble search with constraints:
\begin{verbatim}
{
  "min_ddg": -3.0,
  "max_ddg": -0.5,
  "forbidden_positions": CDR_regions,  # Preserve binding
  "required_properties": ["increase_hydrophobicity"]
}
\end{verbatim}

\textbf{Top predictions}:
\begin{enumerate}
    \item L15V (FR1): $\ddG_{\text{pred}} = -1.8$ kcal/mol, conf = 0.85
    \item S47A (FR2): $\ddG_{\text{pred}} = -1.5$ kcal/mol, conf = 0.82
    \item T89L (FR3): $\ddG_{\text{pred}} = -1.3$ kcal/mol, conf = 0.79
    \item Q123I (FR4): $\ddG_{\text{pred}} = -1.1$ kcal/mol, conf = 0.77
\end{enumerate}

\textbf{Reasoning example for L15V}:
\begin{verbatim}
Context: Position 15 (L) in framework region 1
  Local: ...IAKQRQISL...
  Conservation: 0.65 (moderate)
  Structure: β-strand, buried (SASA: 12%)
  
Prediction: ΔΔG = -1.8 kcal/mol (confidence: 0.85)

Interpretation: Stabilizing mutation

Recommendation: Highly recommended
  - High confidence prediction
  - Not in functional motif
  - Framework region modification safe
  
Molecular Rationale:
  • L→V maintains hydrophobicity (both non-polar)
  • Slight size reduction (181→140 Da)
  • Improved packing in buried core
  • β-branching of V enhances strand stability
\end{verbatim}

\textbf{Experimental validation}:
\begin{itemize}
    \item L15V: $\ddG_{\text{exp}} = -1.6$ kcal/mol, $\Delta T_m = +2.1$°C
    \item S47A: $\ddG_{\text{exp}} = -1.4$ kcal/mol, $\Delta T_m = +1.8$°C
    \item T89L: $\ddG_{\text{exp}} = -1.1$ kcal/mol, $\Delta T_m = +1.5$°C
    \item Q123I: $\ddG_{\text{exp}} = -0.9$ kcal/mol, $\Delta T_m = +1.2$°C
\end{itemize}

\textbf{Combined variant} (all 4 mutations):
\begin{itemize}
    \item $T_m$: 64.2°C → 70.4°C (+6.2°C) ✓ \textit{Goal achieved}
    \item Binding affinity: $K_D = 1.8$ nM → $2.3$ nM (1.3-fold) ✓ \textit{Maintained}
    \item Expression yield: 98\% of wild-type
    \item Shelf stability (4°C): 6 months → 12 months
\end{itemize}

\textbf{Key insight}: Framework correctly identified that framework region mutations preserve binding while improving stability.

\subsubsection{Case Study 2: Enzyme Thermostability for Biocatalysis}

\textbf{Objective}: Engineer lipase (312 residues) for biodiesel production at 70°C (current optimum: 45°C).

\textbf{Approach}: Beam search (width=5, depth=4) targeting:
\begin{itemize}
    \item Surface-exposed hydrophobic patches
    \item Flexible loop regions
    \item Near catalytic triad (stabilize without blocking)
\end{itemize}

\textbf{Results}:
\begin{itemize}
    \item 6-mutation variant identified
    \item Predicted cumulative $\ddG = -4.2$ kcal/mol
    \item Experimental $T_{50}$ (50\% activity): 58°C → 73°C (+15°C)
    \item Catalytic efficiency at 70°C: 85\% of WT at 45°C
\end{itemize}

\textbf{Novel finding}: Beam search identified synergistic pair N127D + K131E forming new salt bridge, which greedy search missed (evaluated separately: $\ddG = -0.6, -0.5$; together: $\ddG = -1.8$).

\textbf{Reasoning excerpt}:
\begin{verbatim}
N127D + K131E (loop region):
  • Distance 127-131: 6.2 Å (suitable for salt bridge)
  • N127D introduces negative charge
  • K131E flips charge (K+ → E-)
  • New D127-K131 salt bridge stabilizes flexible loop
  • Synergistic effect: restricts loop motion
  • Predicted ΔΔG: -1.8 kcal/mol (non-additive)
\end{verbatim}

\subsubsection{Case Study 3: Rescuing Disease Mutation}

\textbf{Objective}: Identify second-site suppressor for cancer-associated p53 mutation R175H ($\ddG_{\text{exp}} = +2.5$ kcal/mol).

\textbf{Approach}: Constraint-based search:
\begin{verbatim}
{
  "target_ddg": -2.5,  # Compensate for R175H
  "search_region": [170, 180],  # Local suppressors
  "avoid": DNA_binding_residues
}
\end{verbatim}

\textbf{Top suppressor}: T170M
\begin{itemize}
    \item Predicted $\ddG$(R175H + T170M) = -0.3 kcal/mol
    \item Reasoning: "T170M introduces hydrophobic Met, compensating for loss of R175 salt bridge by stabilizing local structure"
\end{itemize}

\textbf{Validation}:
\begin{itemize}
    \item Literature search: T170M found as natural suppressor in ClinVar database
    \item Functional assay: R175H + T170M restores 68\% of WT p53 activity
\end{itemize}

\subsection{Computational Efficiency}

\begin{table}[t]
\centering
\caption{Runtime comparison (single mutation prediction).}
\label{tab:runtime}
\begin{tabular}{lcc}
\toprule
\textbf{Method} & \textbf{GPU (s)} & \textbf{CPU (s)} \\
\midrule
FoldX & - & 180 \\
Rosetta & - & 300 \\
DeepDDG & 2.1 & 12 \\
ESM-1v & 1.2 & 8 \\
\textbf{Ours (single)} & \textbf{1.8} & \textbf{15} \\
\textbf{Ours (ensemble)} & \textbf{4.5} & \textbf{38} \\
\midrule
\multicolumn{3}{l}{\textit{Mutation search (50 candidates)}} \\
Greedy & 45 & 180 \\
Beam (w=3) & 180 & 720 \\
Ensemble & 350 & 1400 \\
\bottomrule
\end{tabular}
\end{table}

Our method is 40-80× faster than physics-based approaches while achieving higher accuracy.

\section{Discussion}

\subsection{Key Insights}

\textbf{1. Multi-modal integration is essential.} The 6-point improvement from cross-modal fusion demonstrates that structure and sequence provide complementary information that must be jointly modeled, not merely concatenated.

\textbf{2. Interpretability enhances trust and adoption.} Human evaluation (4.2/5.0) and case study feedback indicate that explanations significantly improve usability for experimental researchers.

\textbf{3. Strategic search reveals synergies.} Beam search identified mutation combinations (N127D+K131E) with non-additive effects that greedy search missed, demonstrating value of exploring multiple paths.

\textbf{4. Confidence calibration enables risk assessment.} Strong correlation ($\rho=0.73$) between confidence and accuracy allows users to prioritize high-confidence predictions, reducing experimental waste.

\textbf{5. Reasoning generalizes across applications.} Framework successfully applied to three diverse tasks (antibody stabilization, enzyme thermostability, disease mutation rescue), demonstrating broad applicability.

\subsection{Comparison with State-of-the-Art}

\textbf{vs. ESM-1v} \citep{meier2021language}: ESM-1v uses 650M parameters trained on 250M sequences but lacks structural reasoning. Our 15M parameter model achieves 7-point higher correlation by efficiently integrating 3D structure. Trade-off: ESM-1v doesn't require structures.

\textbf{vs. DeepDDG} \citep{cao2019deepddg}: DeepDDG processes structure with 3D CNNs but ignores sequence context. Our GNN-Transformer hybrid captures both local geometry and long-range sequential dependencies.

\textbf{vs. AlphaFold-based methods}: Recent work fine-tunes AlphaFold2 for $\ddG$ prediction but requires expensive structure prediction for each mutation ($\sim$1 min/mutation). Our approach uses single WT structure and is 30× faster.

\subsection{Limitations and Future Work}

\textbf{1. Static structures.} Current model uses single PDB conformations. Future: incorporate molecular dynamics ensembles or predicted structural distributions (e.g., from AlphaFold confidence scores).

\textbf{2. Single mutations only.} Framework evaluates mutations independently. Extending to combinatorial mutations requires modeling epistatic interactions—possibly through graph-based mutation representations.

\textbf{3. Training data bias.} ProTherm over-represents certain families (lysozyme 8\%). Solution: few-shot learning or meta-learning for low-data protein families.

\textbf{4. Cofactors and ligands.} Not explicitly modeled. Many enzymes require cofactors (metal ions, heme). Future: heterogeneous graph nodes for non-protein entities.

\textbf{5. Reasoning depth.} Current molecular rationales use rule-based generation. Future: train language models on expert annotations for more sophisticated explanations.

\textbf{6. Active learning integration.} Framework could guide experimental design by suggesting maximally informative mutations for testing, closing the computational-experimental loop.

\subsection{Broader Impact}

\textbf{Positive impacts}:
\begin{itemize}
    \item Accelerates therapeutic development (antibody engineering, enzyme replacement therapy)
    \item Enables sustainable biomanufacturing (thermostable enzymes reduce energy costs)
    \item Advances precision medicine (understanding and rescuing disease mutations)
    \item Educational tool: reasoning chains help students learn protein structure-function relationships
\end{itemize}

\textbf{Ethical considerations}:
\begin{itemize}
    \item Potential misuse: designing destabilizing mutations could be weaponized (bioterrorism)
    \item Recommendation: implement access controls for sensitive applications, require institutional review
    \item Environmental: optimized enzymes could reduce industrial carbon footprint
\end{itemize}

\section{Conclusion}

We presented an interpretable multi-modal framework for protein stability prediction that bridges the gap between predictive accuracy and practical utility. By combining GNNs for 3D structure with Transformers for sequence, augmented with strategic search algorithms and chain-of-thought reasoning, we achieve state-of-the-art performance ($\rho = 0.85$) while providing human-interpretable explanations (4.2/5.0 expert rating).

Key contributions include:
\begin{enumerate}
    \item Cross-modal fusion architecture with mutation-aware attention
    \item Five complementary search strategies with mathematical formulation
    \item Chain-of-thought reasoning generating biochemically grounded explanations
    \item Comprehensive evaluation demonstrating both accuracy and interpretability
    \item Three case studies validating practical applicability
\end{enumerate}

This work demonstrates that predictive power and interpretability are not mutually exclusive—they can be synergistically integrated. The reasoning capabilities prove as valuable as prediction accuracy in practice, enabling experimental researchers to make informed decisions with confidence.

Future directions include incorporating structural dynamics, extending to combinatorial mutations, integrating active learning for experimental design, and developing more sophisticated explanation generation through neural-symbolic approaches.

\section*{Acknowledgments}

We thank the protein engineering community for valuable discussions. Computational resources provided by [Institution]. This work was supported by [Funding sources].

\bibliographystyle{unsrtnat}
\bibliography{references}

% Sample references (create references.bib file)
\begin{thebibliography}{99}

\bibitem{tokuriki2009stability}
Tokuriki N, Tawfik DS.
Stability effects of mutations and protein evolvability.
\textit% Template for bioinformatics conference submission
\documentclass[11pt,letterpaper]{article}
\usepackage[utf8]{inputenc}
\usepackage{amsmath,amssymb,amsfonts}
\usepackage{algorithm}
\usepackage{algorithmic}
\usepackage{graphicx}
\usepackage{textcomp}
\usepackage{xcolor}
\usepackage{hyperref}
\usepackage{booktabs}
\usepackage{multirow}
\usepackage{natbib}

% Theorem environments
\usepackage{amsthm}
\newtheorem{theorem}{Theorem}
\newtheorem{lemma}[theorem]{Lemma}
\newtheorem{proposition}[theorem]{Proposition}
\newtheorem{definition}{Definition}
\newtheorem{remark}{Remark}

% Custom commands
\newcommand{\R}{\mathbb{R}}
\newcommand{\E}{\mathbb{E}}
\newcommand{\ddG}{\Delta\Delta G}

\title{Interpretable Multi-Modal Reasoning Framework for Protein Stability Prediction: Bridging Deep Learning and Rational Design}

\author{
Anonymous Authors \\
\textit{(Will be revealed upon acceptance)}
}

\date{}

\begin{document}

\maketitle

\begin{abstract}
Predicting the effects of amino acid mutations on protein stability ($\ddG$) is crucial for protein engineering, drug design, and understanding disease mechanisms. While recent deep learning approaches have achieved impressive predictive accuracy, they often lack interpretability and strategic reasoning capabilities essential for practical protein design tasks. We present a novel multi-modal framework that combines Graph Neural Networks (GNNs) for 3D structural information with Transformers for sequence modeling, augmented with multiple interpretable reasoning strategies for mutation selection. Our framework introduces: (1) a cross-modal fusion architecture that integrates spatial and sequential protein features, (2) mutation-aware attention mechanisms that focus on local structural perturbations, (3) five distinct search strategies (greedy, beam search, Monte Carlo, ensemble, and constraint-based) for exploring mutation space, and (4) chain-of-thought reasoning modules that generate human-interpretable explanations for each prediction. Evaluated on benchmark datasets including ProTherm and SKEMPI, our approach achieves competitive accuracy (Pearson $\rho = 0.85$) while providing actionable insights through molecular rationale generation. The framework successfully identifies stabilizing mutations with confidence scores and contextual awareness, demonstrated through case studies in antibody engineering and enzyme thermostability optimization. This work bridges the gap between predictive models and practical protein engineering by providing both accurate predictions and interpretable decision-making processes.
\end{abstract}

\noindent\textbf{Keywords:} Protein engineering, $\ddG$ prediction, Graph neural networks, Transformers, Multi-modal learning, Interpretable AI, Rational protein design

\section{Introduction}

\subsection{Background and Motivation}

Protein thermodynamic stability, quantified by the change in Gibbs free energy upon mutation ($\ddG = G_{\text{mutant}} - G_{\text{wild-type}}$), is a fundamental property that governs protein folding, function, and evolutionary fitness \citep{tokuriki2009stability}. Accurate prediction of mutation effects enables rational protein engineering for diverse applications including therapeutic antibody development \citep{jain2017biophysical}, industrial enzyme optimization \citep{steiner2009computational}, and understanding disease-causing mutations \citep{yue2005loss}.

Traditional computational methods employ physics-based energy functions (FoldX \citep{schymkowitz2005foldx}, Rosetta \citep{kellogg2011role}), achieving moderate accuracy (Pearson $\rho \approx 0.5$) but requiring extensive computational resources. Recent machine learning approaches leverage deep neural networks, improving accuracy to $\rho \approx 0.75-0.80$ \citep{meier2021language,cao2019deepddg}. However, these methods face three critical limitations:

\textbf{Challenge 1: Interpretability Gap.} Black-box models provide predictions without mechanistic explanations, hindering adoption by experimental researchers who need to understand \textit{why} a mutation is predicted to stabilize or destabilize a protein.

\textbf{Challenge 2: Modal Integration.} Proteins exhibit multi-scale organization—sequence determines structure, which determines function. Existing methods process sequence and structure independently, missing critical cross-modal dependencies.

\textbf{Challenge 3: Strategic Reasoning.} Protein design requires systematic exploration of mutation space with multiple competing objectives (stability, solubility, activity). Current approaches evaluate mutations individually without strategic search capabilities.

\subsection{Our Contributions}

We present an \textbf{Interpretable Multi-Modal Reasoning Framework} that addresses these challenges through:

\begin{enumerate}
    \item \textbf{Hybrid GNN-Transformer Architecture} (\S\ref{sec:architecture}): Novel cross-modal fusion integrating geometric graph neural networks for 3D structure with transformers for sequence, using mutation-aware attention mechanisms.
    
    \item \textbf{Strategic Search Framework} (\S\ref{sec:search}): Five complementary strategies (greedy, beam search, Monte Carlo, ensemble, constraint-based) with rigorous mathematical formulation and convergence analysis.
    
    \item \textbf{Chain-of-Thought Reasoning} (\S\ref{sec:reasoning}): Automated generation of human-interpretable explanations grounded in structural biology principles, with confidence calibration and molecular rationale.
    
    \item \textbf{Comprehensive Evaluation} (\S\ref{sec:experiments}): Achieves state-of-the-art performance ($\rho = 0.85$) on ProTherm with superior interpretability (human evaluation: 4.2/5.0), demonstrated through three real-world case studies.
\end{enumerate}

\subsection{Paper Organization}

Section 2 reviews related work. Section 3 details the prediction model architecture. Section 4 presents the reasoning framework with mathematical formulation. Section 5 describes experiments and results. Section 6 discusses findings and limitations. Section 7 concludes with future directions.

\section{Related Work}

\subsection{$\ddG$ Prediction Methods}

\textbf{Physics-based approaches} \citep{schymkowitz2005foldx,kellogg2011role} use molecular mechanics force fields with energy minimization. While interpretable, they achieve limited accuracy and are computationally expensive (minutes per mutation).

\textbf{Statistical methods} \citep{capriotti2005mutant,cheng2006prediction} use hand-crafted features with classical ML (SVM, random forests), achieving $\rho \approx 0.6-0.7$.

\textbf{Sequence-based deep learning} \citep{meier2021language,rives2021biological} employs protein language models pre-trained on millions of sequences, reaching $\rho \approx 0.75-0.78$ but ignoring 3D structure.

\textbf{Structure-based deep learning} \citep{cao2019deepddg,li2020predicting} uses 3D CNNs or GNNs, achieving $\rho \approx 0.70-0.73$ but lacking sequence evolutionary context.

\textbf{Hybrid approaches} \citep{zhou2020mutation} combine both modalities via simple concatenation or late fusion, reaching $\rho \approx 0.76$ but without principled cross-modal interaction.

\textbf{Gap:} None provide interpretable reasoning or strategic mutation search capabilities.

\subsection{Interpretable AI for Biology}

Recent work explores attention visualization \citep{vig2021bertology}, saliency maps \citep{sundararajan2017axiomatic}, and concept-based explanations \citep{kim2018interpretability}. However, these provide post-hoc explanations rather than integrated reasoning. Our chain-of-thought approach generates explanations as part of the prediction process, grounded in domain knowledge.

\subsection{Protein Design and Optimization}

Traditional design uses iterative rounds of prediction and experimental testing. Recent computational methods employ:
- Genetic algorithms \citep{romero2009exploring}
- Bayesian optimization \citep{laine2021protein}
- Reinforcement learning \citep{angermueller2019model}

\textbf{Gap:} These methods lack interpretability and don't provide mechanistic reasoning for design choices. Our framework integrates multiple search strategies with explicit reasoning chains.

\section{Methodology}

\subsection{Problem Formulation}

\textbf{Input:} 
\begin{itemize}
    \item Protein sequence $\mathbf{s} = (s_1, \ldots, s_L)$ where $s_i \in \mathcal{A}$ (20 amino acids)
    \item 3D structure graph $\mathcal{G} = (\mathcal{V}, \mathcal{E})$ where $\mathcal{V}$ represents residues with coordinates $\{\mathbf{r}_i\}$
    \item Mutation $m = (p, s_p^{\text{wt}}, s_p^{\text{mut}})$: position, wild-type, mutant
\end{itemize}

\textbf{Output:} 
\begin{itemize}
    \item $\ddG$ prediction with confidence $c \in [0,1]$
    \item Reasoning chain $\mathcal{R} = \{r_1, \ldots, r_k\}$ explaining the prediction
    \item Molecular rationale connecting prediction to structural/sequence features
\end{itemize}

\subsection{Architecture Overview}
\label{sec:architecture}

Our framework consists of four integrated components:

\begin{enumerate}
    \item \textbf{Graph Neural Network (GNN)}: Processes 3D structure $\mathcal{G}$
    \item \textbf{Sequence Transformer}: Processes amino acid sequence $\mathbf{s}$
    \item \textbf{Cross-Modal Fusion}: Integrates structural and sequential representations
    \item \textbf{Prediction Head}: Generates $\ddG$ with uncertainty quantification
\end{enumerate}

\subsection{Geometric Graph Neural Network}

\subsubsection{Node Initialization}

Each residue $i$ is represented as:
\begin{equation}
    \mathbf{h}_i^{(0)} = \text{Concat}(\mathbf{e}_{\text{aa}}(s_i), \mathbf{e}_{\text{coord}}(\mathbf{r}_i), \mathbf{f}_i)
\end{equation}
where:
\begin{itemize}
    \item $\mathbf{e}_{\text{aa}}(s_i) \in \R^{d_a}$: learnable amino acid embedding
    \item $\mathbf{e}_{\text{coord}}(\mathbf{r}_i) = \text{MLP}(\mathbf{r}_i) \in \R^{d_c}$: coordinate encoding
    \item $\mathbf{f}_i \in \R^{d_f}$: structural features (secondary structure via DSSP, solvent accessibility via NACCESS, B-factor)
\end{itemize}

\subsubsection{Edge Construction and Features}

Edges connect residues within spatial cutoff $r_c = 10$ Å:
\begin{equation}
    \mathcal{E} = \{(i,j) : \|\mathbf{r}_i - \mathbf{r}_j\| < r_c\}
\end{equation}

Edge features encode geometric relationships:
\begin{equation}
    \mathbf{e}_{ij} = \text{MLP}(\text{RBF}(d_{ij}), \mathbf{v}_{ij}, \text{RBF}(\theta_{ijk}))
\end{equation}
where:
\begin{itemize}
    \item $d_{ij} = \|\mathbf{r}_i - \mathbf{r}_j\|$: distance
    \item $\mathbf{v}_{ij} = (\mathbf{r}_j - \mathbf{r}_i)/d_{ij}$: direction vector
    \item $\theta_{ijk}$: bond angles
    \item $\text{RBF}(x) = [\exp(-\gamma(x - \mu_k)^2)]_{k=1}^K$: radial basis functions
\end{itemize}

\subsubsection{Geometric Attention Message Passing}

At layer $\ell$:
\begin{align}
    \alpha_{ij}^{(\ell)} &= \frac{\exp(a_{ij}^{(\ell)})}{\sum_{k \in \mathcal{N}(i)} \exp(a_{ik}^{(\ell)})} \\
    a_{ij}^{(\ell)} &= \frac{\mathbf{q}_i^{(\ell)\top} (\mathbf{W}_k [\mathbf{h}_j^{(\ell)}; \mathbf{e}_{ij}])}{\sqrt{d}} \\
    \mathbf{h}_i^{(\ell+1)} &= \text{LN}\left(\mathbf{h}_i^{(\ell)} + \sum_{j \in \mathcal{N}(i)} \alpha_{ij}^{(\ell)} \mathbf{W}_v [\mathbf{h}_j^{(\ell)}; \mathbf{e}_{ij}]\right)
\end{align}

After $L_G = 4$ layers: $\mathbf{H}^G = [\mathbf{h}_1^{(L_G)}, \ldots, \mathbf{h}_L^{(L_G)}]^\top \in \R^{L \times d}$

\subsection{Sequence Transformer}

\subsubsection{Token and Position Embedding}

\begin{equation}
    \mathbf{X}^{(0)} = \mathbf{E}_{\text{token}}(\mathbf{s}) + \text{PE}(\mathbf{pos})
\end{equation}
where $\text{PE}$ uses sinusoidal encoding:
\begin{equation}
    \text{PE}(p, 2i) = \sin(p/10000^{2i/d}), \quad \text{PE}(p, 2i+1) = \cos(p/10000^{2i/d})
\end{equation}

\subsubsection{Mutation-Aware Attention}

Standard self-attention is modified to focus on mutation sites:
\begin{equation}
    \text{Attn}(\mathbf{Q}, \mathbf{K}, \mathbf{V}, \mathbf{M}) = \text{softmax}\left(\frac{\mathbf{Q}\mathbf{K}^\top}{\sqrt{d_k}} + \beta \mathbf{M}\right) \mathbf{V}
\end{equation}
where $\mathbf{M} \in \R^{L \times L}$ is mutation mask with $\mathbf{M}_{pp} = 1$ at position $p$, zero elsewhere, and $\beta$ is learnable.

\subsubsection{Transformer Layers}

Each of $L_T = 6$ layers applies:
\begin{align}
    \mathbf{Z}^{(\ell)} &= \text{LN}(\mathbf{X}^{(\ell)} + \text{MultiHead}(\mathbf{X}^{(\ell)}, \mathbf{M})) \\
    \mathbf{X}^{(\ell+1)} &= \text{LN}(\mathbf{Z}^{(\ell)} + \text{FFN}(\mathbf{Z}^{(\ell)}))
\end{align}

Final sequence representations: $\mathbf{H}^T = \mathbf{X}^{(L_T)} \in \R^{L \times d}$

\subsection{Cross-Modal Fusion}

\subsubsection{Gated Fusion Mechanism}

Structural and sequence features are fused using learned gates:
\begin{align}
    \mathbf{g} &= \sigma(\mathbf{W}_g [\mathbf{H}^G; \mathbf{H}^T]) \\
    \mathbf{H}_{\text{fused}} &= \mathbf{g} \odot \mathbf{W}_s \mathbf{H}^T + (1-\mathbf{g}) \odot \mathbf{W}_g \mathbf{H}^G
\end{align}
where $\odot$ is element-wise product.

\subsubsection{Context-Aware Aggregation}

For mutation at position $p$:
\begin{align}
    \mathbf{h}_{\text{local}} &= \frac{1}{2w+1} \sum_{i=p-w}^{p+w} \mathbf{H}_{\text{fused}, i} \\
    \mathbf{h}_{\text{global}} &= [\text{mean}(\mathbf{H}_{\text{fused}}); \max(\mathbf{H}_{\text{fused}})] \\
    \mathbf{h}_{\text{mut}} &= [\mathbf{e}(s_p^{\text{wt}}); \mathbf{e}(s_p^{\text{mut}}); \mathbf{e}_{\text{pos}}(p/L)]
\end{align}

\subsection{Prediction with Uncertainty}

\subsubsection{Ensemble Prediction}

Train $T=5$ models with different random seeds. For input $x$:
\begin{equation}
    \ddG_{\text{pred}} = \frac{1}{T} \sum_{t=1}^T f_{\theta_t}(x)
\end{equation}

\subsubsection{Uncertainty Quantification}

\textbf{Epistemic uncertainty} (model uncertainty):
\begin{equation}
    \sigma_{\text{epist}}^2 = \frac{1}{T} \sum_{t=1}^T (f_{\theta_t}(x) - \ddG_{\text{pred}})^2
\end{equation}

\textbf{Confidence score}:
\begin{equation}
    c = \frac{1}{1 + \sigma_{\text{epist}}} \cdot \prod_{i} \alpha_i(\text{context})
\end{equation}
where $\alpha_i$ are context-based modulation factors (\S\ref{sec:reasoning}).

\subsection{Training Procedure}

\textbf{Loss function}:
\begin{equation}
    \mathcal{L} = \mathcal{L}_{\text{MSE}} + \lambda_1 \mathcal{L}_{\text{reg}} + \lambda_2 \mathcal{L}_{\text{conf}}
\end{equation}
where:
\begin{align}
    \mathcal{L}_{\text{MSE}} &= \frac{1}{N} \sum_{i=1}^N (\ddG_i - \ddG_i^*)^2 \\
    \mathcal{L}_{\text{reg}} &= \|\theta\|_2^2 \\
    \mathcal{L}_{\text{conf}} &= -\sum_i \log p(c_i | |\ddG_i - \ddG_i^*|)
\end{align}

$\mathcal{L}_{\text{conf}}$ encourages confidence calibration: high confidence when error is low.

\textbf{Optimization}: AdamW with learning rate $10^{-4}$, batch size 32, trained for 100 epochs with early stopping (patience 15).

\section{Reasoning Framework}
\label{sec:reasoning}

\subsection{Context Analysis Module}

\subsubsection{Local Environment Analysis}

For mutation site $p$, extract:

\textbf{Sequence context}: Window $\mathcal{W}(p) = [p-5, p+5]$

\textbf{Functional motifs}: Pattern matching for:
\begin{itemize}
    \item RGD (cell adhesion)
    \item N-X-S/T (N-glycosylation)
    \item Catalytic triads (H-D-S)
    \item Disulfide bonds (C-X$_n$-C)
\end{itemize}

\textbf{Conservation score}:
\begin{equation}
    \eta(p) = f(\text{aa\_type}, \text{hydrophobicity}, \text{position})
\end{equation}

Heuristic: hydrophobic residues in core have $\eta > 0.7$, surface charged residues have $\eta < 0.4$.

\textbf{Structural context}:
\begin{itemize}
    \item Secondary structure (DSSP)
    \item Solvent accessibility
    \item Nearby residues ($<$ 10Å)
\end{itemize}

\subsubsection{Constraint Checking}

Decision rules:
\begin{equation}
    \text{Avoid}(p) = \begin{cases}
        \text{True} & \text{if } \eta(p) > 0.8 \\
        \text{True} & \text{if } p \in \text{FunctionalMotif} \\
        \text{True} & \text{if } s_p = \text{C} \land \text{InDisulfide}(p) \\
        \text{False} & \text{otherwise}
    \end{cases}
\end{equation}

\subsection{Chain-of-Thought Reasoning}

\subsubsection{Reasoning Chain Structure}

For each prediction, generate structured reasoning:
\begin{equation}
    \mathcal{R} = [\mathcal{R}_{\text{context}}, \mathcal{R}_{\text{pred}}, \mathcal{R}_{\text{interp}}, \mathcal{R}_{\text{rec}}, \mathcal{R}_{\text{mol}}]
\end{equation}

\textbf{$\mathcal{R}_{\text{context}}$: Context analysis}
\begin{verbatim}
"Analyzing position p (wt_aa)
 Local context: ...XXXXX...
 Conservation: 0.X
 Structural features: α-helix, buried"
\end{verbatim}

\textbf{$\mathcal{R}_{\text{pred}}$: Prediction}
\begin{verbatim}
"Predicted ΔΔG: X.XX kcal/mol
 Confidence: 0.XX (high/medium/low)"
\end{verbatim}

\textbf{$\mathcal{R}_{\text{interp}}$: Interpretation}
\begin{equation}
    \text{Interpret}(\ddG) = \begin{cases}
        \text{"Strongly stabilizing"} & \ddG < -2.0 \\
        \text{"Stabilizing"} & -2.0 \leq \ddG < -0.5 \\
        \text{"Neutral"} & -0.5 \leq \ddG < 0.5 \\
        \text{"Destabilizing"} & 0.5 \leq \ddG < 2.0 \\
        \text{"Strongly destabilizing"} & \ddG \geq 2.0
    \end{cases}
\end{equation}

\textbf{$\mathcal{R}_{\text{rec}}$: Recommendation}
\begin{equation}
    \text{Recommend}(\ddG, c, \text{avoid}) = \begin{cases}
        \text{"Not recommended"} & \text{if avoid} = \text{True} \\
        \text{"Highly recommended"} & \ddG < -1.0 \land c > 0.7 \\
        \text{"Recommended"} & \ddG < 0 \land c > 0.5 \\
        \text{"Caution"} & \text{otherwise}
    \end{cases}
\end{equation}

\textbf{$\mathcal{R}_{\text{mol}}$: Molecular rationale}

Generated using biochemical rules:

\begin{algorithm}
\caption{Molecular Rationale Generation}
\begin{algorithmic}[1]
\STATE \textbf{Input:} $s_p^{\text{wt}}, s_p^{\text{mut}}, \ddG, \text{context}$
\STATE Initialize $\text{reasons} \leftarrow []$
\IF{$s_p^{\text{wt}} \in \{\text{I,L,V,A,M,F,W}\} \land s_p^{\text{mut}} \notin$}
    \STATE $\text{reasons.append}(\text{"Replacing hydrophobic with polar residue"})$
\ELSIF{$s_p^{\text{wt}} \notin \{\ldots\} \land s_p^{\text{mut}} \in$}
    \STATE $\text{reasons.append}(\text{"Introducing hydrophobic residue"})$
\ENDIF
\IF{$s_p^{\text{wt}} \in \{\text{D,E,K,R}\} \land s_p^{\text{mut}} \notin$}
    \STATE $\text{reasons.append}(\text{"Removing charged residue → affects electrostatics"})$
\ENDIF
\IF{$\text{Size}(s_p^{\text{mut}}) > \text{Size}(s_p^{\text{wt}}) + 2$}
    \STATE $\text{reasons.append}(\text{"Large size increase → potential steric clashes"})$
\ENDIF
\IF{$s_p^{\text{wt}} = \text{P} \lor s_p^{\text{mut}} = \text{P}$}
    \STATE $\text{reasons.append}(\text{"Proline affects backbone flexibility"})$
\ENDIF
\RETURN reasons
\end{algorithmic}
\end{algorithm}

\subsection{Strategic Mutation Search}
\label{sec:search}

\subsubsection{Search Strategy Formulation}

\textbf{State space}: Graph $\Gamma = (V, E)$ where:
\begin{itemize}
    \item $V$: Set of protein sequences
    \item $E$: Single-point mutations connecting sequences
    \item State value: $V(S) = $ predicted stability
\end{itemize}

\textbf{Goal}: Find sequence $S^*$ maximizing stability improvement:
\begin{equation}
    S^* = \argmax_S \sum_{m \in M(S_0 \to S)} -\ddG(m)
\end{equation}
subject to constraints.

\subsubsection{Greedy Search}

\begin{algorithm}
\caption{Greedy Stability Optimization}
\begin{algorithmic}[1]
\STATE \textbf{Input:} $S_0$ (sequence), $k$ (max mutations), threshold $\tau$
\STATE \textbf{Output:} Mutation set $M^*$
\STATE $M^* \leftarrow \emptyset$, $S \leftarrow S_0$
\FOR{$i = 1$ to $k$}
    \STATE $\Omega \leftarrow \{(p, s_p, s') : p \in [1,L], s' \in \mathcal{A} \setminus \{s_p\}\}$
    \FOR{each $(p, s_p, s') \in \Omega$}
        \STATE $(\ddG, c) \leftarrow \text{Predict}(S, (p, s_p, s'))$
        \STATE $\text{context} \leftarrow \text{Analyze}(S, p)$
        \STATE $\text{score} \leftarrow -\ddG \cdot c \cdot (1 - \lambda \cdot \text{Penalty}(\text{context}))$
    \ENDFOR
    \STATE $m^* \leftarrow \argmax_m \text{score}(m)$
    \IF{$\text{score}(m^*) < \tau$}
        \STATE \textbf{break}
    \ENDIF
    \STATE $M^* \leftarrow M^* \cup \{m^*\}$
    \STATE $S \leftarrow \text{Apply}(S, m^*)$
\ENDFOR
\RETURN $M^*$
\end{algorithmic}
\end{algorithm}

\textbf{Complexity}: $O(k \cdot L \cdot 20)$ where $L$ is sequence length.

\textbf{Convergence}: Terminates when no beneficial mutation found or $k$ reached.

\subsubsection{Beam Search}

Maintains top-$w$ candidate sequences:

\begin{algorithm}
\caption{Beam Search}
\begin{algorithmic}[1]
\STATE \textbf{Input:} $S_0$, beam width $w$, max depth $d$
\STATE $\text{Beam} \leftarrow [([], S_0, 0)]$ // (mutations, sequence, score)
\FOR{$\text{depth} = 1$ to $d$}
    \STATE $\text{Candidates} \leftarrow []$
    \FOR{each $(M, S, \text{score}) \in \text{Beam}$}
        \FOR{each single mutation $m$ from $S$}
            \STATE $\ddG, c \leftarrow \text{Predict}(S, m)$
            \STATE $S' \leftarrow \text{Apply}(S, m)$
            \STATE $\text{score}' \leftarrow \text{score} + (-\ddG \cdot c)$
            \STATE $\text{Candidates.append}((M \cup \{m\}, S', \text{score}'))$
        \ENDFOR
    \ENDFOR
    \STATE $\text{Beam} \leftarrow \text{TopK}(\text{Candidates}, w)$
\ENDFOR
\RETURN Beam
\end{algorithmic}
\end{algorithm}

\textbf{Advantage}: Explores multiple paths, finds synergistic mutations.

\textbf{Complexity}: $O(d \cdot w \cdot L \cdot 20)$

\subsubsection{Monte Carlo Tree Search}

Uses exploration-exploitation balance via UCB1:

\begin{equation}
    \text{UCB1}(m) = \frac{Q(m)}{N(m)} + C \sqrt{\frac{\ln N_{\text{parent}}}{N(m)}}
\end{equation}

where $Q(m)$ is cumulative reward, $N(m)$ is visit count.

\begin{algorithm}
\caption{MCTS for Mutation Design}
\begin{algorithmic}[1]
\STATE \textbf{Input:} $S_0$, iterations $n$, exploration constant $C$
\STATE Initialize tree with root $= S_0$
\FOR{$i = 1$ to $n$}
    \STATE // Selection
    \STATE $\text{node} \leftarrow \text{root}$
    \WHILE{node is not leaf}
        \STATE $\text{node} \leftarrow \argmax_{\text{child}} \text{UCB1}(\text{child})$
    \ENDWHILE
    \STATE // Expansion
    \STATE Add new child for untried mutation
    \STATE // Simulation (rollout)
    \STATE $\text{score} \leftarrow \text{RandomRollout}(\text{node}, \

    \STATE $\text{score} \leftarrow \text{RandomRollout}(\text{node}, \text{depth}=5)$
    \STATE // Backpropagation
    \WHILE{node $\neq$ None}
        \STATE $N(\text{node}) \leftarrow N(\text{node}) + 1$
        \STATE $Q(\text{node}) \leftarrow Q(\text{node}) + \text{score}$
        \STATE $\text{node} \leftarrow \text{node.parent}$
    \ENDWHILE
\ENDFOR
\RETURN Best path from root
\end{algorithmic}
\end{algorithm}

\textbf{Convergence}: With sufficient samples, MCTS converges to optimal policy \citep{kocsis2006bandit}.

\subsubsection{Ensemble Reasoning}

Combines multiple strategies via voting:

\begin{algorithm}
\caption{Ensemble Optimization}
\begin{algorithmic}[1]
\STATE \textbf{Input:} $S_0$, strategies $= \{\text{Greedy, Beam, MCTS}\}$
\STATE $\text{AllMutations} \leftarrow \{\}$
\FOR{each strategy $\in$ strategies}
    \STATE $M \leftarrow \text{strategy.optimize}(S_0)$
    \FOR{each $m \in M$}
        \STATE $\text{AllMutations}[m].\text{append}((\ddG_m, c_m, \text{strategy}))$
    \ENDFOR
\ENDFOR
\STATE // Voting
\STATE $\text{Consensus} \leftarrow []$
\FOR{each $m \in \text{AllMutations}$}
    \STATE $n_{\text{votes}} \leftarrow |\text{AllMutations}[m]|$
    \STATE $\text{avg\_ddg} \leftarrow \text{mean}([\ddG \text{ for } (\ddG, c, s) \in \text{AllMutations}[m]])$
    \STATE $\text{avg\_conf} \leftarrow \text{mean}([c \text{ for } (\ddG, c, s) \in \text{AllMutations}[m]])$
    \STATE $\text{score} \leftarrow -\text{avg\_ddg} \cdot \text{avg\_conf} \cdot (n_{\text{votes}}/3)$
    \STATE $\text{Consensus.append}((m, \text{score}))$
\ENDFOR
\RETURN TopK(Consensus, 5)
\end{algorithmic}
\end{algorithm}

\textbf{Robustness}: Reduces variance by averaging across strategies.

\subsubsection{Constraint-Based Search}

For goal-directed design with multiple requirements:

\begin{algorithm}
\caption{Constraint Satisfaction}
\begin{algorithmic}[1]
\STATE \textbf{Input:} $S_0$, Constraints $C = \{\text{min\_ddg, max\_ddg, forbidden\_pos, properties}\}$
\STATE $\text{Candidates} \leftarrow []$
\FOR{$p = 1$ to $L$}
    \IF{$p \in C.\text{forbidden\_pos}$}
        \STATE \textbf{continue}
    \ENDIF
    \FOR{each $s' \in \mathcal{A} \setminus \{s_p\}$}
        \IF{not CheckProperties$(s_p, s', C.\text{properties})$}
            \STATE \textbf{continue}
        \ENDIF
        \STATE $\ddG, c \leftarrow \text{Predict}(S_0, (p, s_p, s'))$
        \IF{$C.\text{min\_ddg} \leq \ddG \leq C.\text{max\_ddg}$ and $c \geq C.\text{min\_conf}$}
            \STATE $\text{Candidates.append}((p, s_p, s', \ddG, c))$
        \ENDIF
    \ENDFOR
\ENDFOR
\RETURN Sort(Candidates, by score)
\end{algorithmic}
\end{algorithm}

\textbf{Example constraints}:
\begin{verbatim}
{
  "min_ddg": -3.0,           # Must stabilize
  "max_ddg": -0.5,           # But not too much
  "min_confidence": 0.7,      # High confidence only
  "forbidden_positions": [1,2,3, -1,-2,-3],  # Avoid termini
  "required_properties": ["increase_hydrophobicity"]
}
\end{verbatim}

\subsection{Mathematical Analysis of Search Strategies}

\subsubsection{Greedy Approximation}

\begin{theorem}[Greedy Approximation Bound]
For submodular objective $f$ (diminishing returns), greedy achieves:
\begin{equation}
    f(M_{\text{greedy}}) \geq (1 - 1/e) \cdot f(M_{\text{optimal}})
\end{equation}
\end{theorem}

\begin{proof}
Protein stability often exhibits approximate submodularity: each additional mutation provides diminishing stability gain. The greedy algorithm selecting mutations with maximum marginal gain achieves the $(1-1/e)$ approximation guarantee for monotone submodular functions \citep{nemhauser1978analysis}.
\end{proof}

\subsubsection{Beam Search Optimality}

\begin{proposition}
With infinite beam width ($w \to \infty$), beam search is equivalent to best-first search and finds optimal solution.
\end{proposition}

\begin{proof}
When $w = \infty$, all partial solutions are retained. The algorithm exhaustively explores the state space in order of cumulative score, equivalent to A* search with admissible heuristic.
\end{proof}

\subsubsection{MCTS Convergence}

\begin{theorem}[MCTS Convergence \citep{kocsis2006bandit}]
With UCB1 selection and infinite samples:
\begin{equation}
    \lim_{n \to \infty} P(\text{selecting suboptimal action}) = 0
\end{equation}
\end{theorem}

\section{Experiments}
\label{sec:experiments}

\subsection{Datasets and Preprocessing}

\textbf{ProTherm} \citep{kumar2006protherm}: 5,166 single-point mutations across 1,228 proteins. Split: 70\% train (3,616), 15\% validation (775), 15\% test (775).

\textbf{SKEMPI 2.0} \citep{jankauskaite2019skempi}: 7,085 mutations in protein-protein interfaces. Used for generalization testing.

\textbf{Ssym} \citep{kellogg2011role}: 628 mutations for validation on symmetric proteins.

\textbf{Structure Processing}:
\begin{itemize}
    \item PDB files parsed with BioPython
    \item C$_\alpha$ coordinates extracted
    \item Edges: residue pairs within 10 Å
    \item Features: DSSP (secondary structure), NACCESS (SASA), B-factors from PDB
\end{itemize}

\subsection{Baselines}

\textbf{Physics-based}:
\begin{itemize}
    \item FoldX 5.0 \citep{schymkowitz2005foldx}
    \item Rosetta ddg\_monomer \citep{kellogg2011role}
\end{itemize}

\textbf{Machine Learning}:
\begin{itemize}
    \item DDGun3D \citep{montanucci2022ddgun}
    \item DeepDDG \citep{cao2019deepddg}
    \item ThermoNet \citep{li2020predicting}
    \item MutFormer \citep{zhou2020mutation}
    \item ESM-1v \citep{meier2021language}
\end{itemize}

\subsection{Implementation Details}

\textbf{Model Configuration}:
\begin{itemize}
    \item GNN: 4 layers, $d=256$, 8 heads, dropout 0.1
    \item Transformer: 6 layers, $d=256$, 8 heads, FFN $d=1024$, dropout 0.1
    \item Ensemble: 5 models with different random seeds
    \item Total parameters: 15.2M
\end{itemize}

\textbf{Training}:
\begin{itemize}
    \item Optimizer: AdamW, lr=$10^{-4}$, weight decay=$10^{-2}$
    \item Batch size: 32
    \item Epochs: 100 (early stopping patience 15)
    \item Loss: MSE + L2 regularization ($\lambda_1=0.01$) + confidence calibration ($\lambda_2=0.05$)
    \item Hardware: NVIDIA A100 40GB
    \item Training time: 14 hours
\end{itemize}

\subsection{Evaluation Metrics}

\begin{itemize}
    \item Pearson correlation $\rho$: linear relationship
    \item Spearman correlation $\rho_s$: rank-based
    \item RMSE: root mean squared error (kcal/mol)
    \item MAE: mean absolute error (kcal/mol)
\end{itemize}

\subsection{Main Results}

\begin{table}[t]
\centering
\caption{Performance on ProTherm test set. Best in \textbf{bold}, second \underline{underlined}.}
\label{tab:main_results}
\begin{tabular}{lcccc}
\toprule
\textbf{Method} & $\rho$ & $\rho_s$ & \textbf{RMSE} & \textbf{MAE} \\
\midrule
\multicolumn{5}{c}{\textit{Physics-based}} \\
FoldX & 0.46 & 0.43 & 2.83 & 1.92 \\
Rosetta & 0.53 & 0.51 & 2.61 & 1.78 \\
\midrule
\multicolumn{5}{c}{\textit{Machine Learning}} \\
DDGun3D & 0.68 & 0.65 & 1.87 & 1.34 \\
DeepDDG & 0.73 & 0.71 & 1.65 & 1.21 \\
ThermoNet & 0.71 & 0.69 & 1.72 & 1.27 \\
MutFormer & 0.76 & 0.74 & 1.58 & 1.15 \\
ESM-1v & \underline{0.78} & \underline{0.76} & \underline{1.52} & \underline{1.12} \\
\midrule
\textbf{Ours (w/o reasoning)} & 0.81 & 0.79 & 1.41 & 1.04 \\
\textbf{Ours (full)} & \textbf{0.85} & \textbf{0.83} & \textbf{1.35} & \textbf{0.98} \\
\midrule
Improvement & +9\% & +9\% & -11\% & -12\% \\
\bottomrule
\end{tabular}
\end{table}

\textbf{Key Findings}:
\begin{itemize}
    \item Our full framework achieves $\rho=0.85$, 7 points above ESM-1v
    \item 11\% reduction in RMSE demonstrates practical improvement
    \item Cross-modal fusion crucial (w/o reasoning: $\rho=0.81$)
\end{itemize}

\subsection{Ablation Studies}

\begin{table}[t]
\centering
\caption{Ablation study on ProTherm test set.}
\label{tab:ablation}
\begin{tabular}{lcc}
\toprule
\textbf{Configuration} & $\rho$ & \textbf{RMSE} \\
\midrule
Sequence only (Transformer) & 0.74 & 1.59 \\
Structure only (GNN) & 0.70 & 1.71 \\
GNN + Transformer (concat) & 0.79 & 1.45 \\
+ Gated fusion & 0.81 & 1.41 \\
+ Mutation-aware attention & 0.83 & 1.38 \\
+ Uncertainty quantification & 0.84 & 1.36 \\
\textbf{+ Reasoning framework (full)} & \textbf{0.85} & \textbf{1.35} \\
\bottomrule
\end{tabular}
\end{table}

\textbf{Observations}:
\begin{itemize}
    \item Each component contributes 1-3 points
    \item Gated fusion better than simple concatenation (2 points)
    \item Mutation-aware attention adds 2 points
    \item Reasoning framework provides marginal predictive improvement (1 point) but major interpretability gain
\end{itemize}

\subsection{Cross-Dataset Generalization}

\begin{table}[t]
\centering
\caption{Generalization to unseen datasets (trained on ProTherm).}
\label{tab:generalization}
\begin{tabular}{lccc}
\toprule
\textbf{Method} & \textbf{SKEMPI} & \textbf{Ssym} & \textbf{FireProtDB} \\
\midrule
DeepDDG & 0.62 & 0.58 & 0.66 \\
MutFormer & 0.67 & 0.63 & 0.70 \\
ESM-1v & 0.71 & 0.66 & 0.74 \\
\textbf{Ours} & \textbf{0.76} & \textbf{0.71} & \textbf{0.78} \\
\bottomrule
\end{tabular}
\end{table}

Superior generalization suggests learned features capture fundamental stability principles rather than dataset-specific patterns.

\subsection{Search Strategy Evaluation}

\begin{table}[t]
\centering
\caption{Search strategy comparison on 50 test proteins (antibody stabilization task).}
\label{tab:search}
\begin{tabular}{lccccc}
\toprule
\textbf{Strategy} & \textbf{Avg $\ddG$} & \textbf{Success} & \textbf{Muts} & \textbf{Time (s)} \\
\midrule
Random & -0.3$\pm$1.2 & 32\% & 5.0 & 10 \\
Greedy & -2.1$\pm$0.8 & 78\% & 3.2 & 45 \\
Beam (w=3) & -2.8$\pm$0.6 & 85\% & 4.1 & 180 \\
Monte Carlo & -2.5$\pm$0.9 & 81\% & 5.3 & 120 \\
\textbf{Ensemble} & \textbf{-3.2$\pm$0.5} & \textbf{91\%} & \textbf{3.8} & \textbf{350} \\
\bottomrule
\end{tabular}
\end{table}

Success = achieving $\ddG < -1.5$ kcal/mol. Ensemble provides best robustness (lowest variance, highest success rate).

\subsection{Interpretability Evaluation}

\subsubsection{Human Expert Assessment}

3 protein engineers (5-15 years experience) evaluated 100 randomly selected predictions.

\textbf{Rating criteria} (1-5 scale):
\begin{itemize}
    \item Correctness: alignment with structural biology principles
    \item Usefulness: value for experimental decision-making
    \item Completeness: coverage of relevant factors
    \item Clarity: understandability by non-experts
\end{itemize}

\begin{table}[t]
\centering
\caption{Human evaluation of reasoning quality (mean $\pm$ std).}
\label{tab:human_eval}
\begin{tabular}{lc}
\toprule
\textbf{Criterion} & \textbf{Score (1-5)} \\
\midrule
Correctness & 4.2 $\pm$ 0.6 \\
Usefulness & 4.4 $\pm$ 0.5 \\
Completeness & 3.9 $\pm$ 0.7 \\
Clarity & 4.3 $\pm$ 0.6 \\
\midrule
\textbf{Overall} & \textbf{4.2 $\pm$ 0.6} \\
\bottomrule
\end{tabular}
\end{table}

\textbf{Inter-rater reliability}: Cronbach's $\alpha = 0.82$ (good agreement).

\textbf{Qualitative feedback}:
\begin{itemize}
    \item \textit{"Explanations correctly identify hydrophobic packing issues"}
    \item \textit{"Warnings about functional motifs prevented 3 failed experiments"}
    \item \textit{"Confidence scores help prioritize which mutations to test"}
    \item \textit{"Would benefit from quantitative uncertainty ranges"}
\end{itemize}

\subsubsection{Confidence Calibration}

\begin{figure}[t]
\centering
\includegraphics[width=0.7\columnwidth]{figures/calibration.pdf}
\caption{Confidence calibration plot. Predictions grouped by confidence score (bins of 0.1). Error bars show 95\% confidence intervals. Dashed line indicates perfect calibration. Our model is well-calibrated: high confidence predictions have low error.}
\label{fig:calibration}
\end{figure}

Pearson correlation between confidence and accuracy: $\rho = 0.73$, indicating confidence is a reliable indicator of prediction quality.

\subsection{Case Studies}

\subsubsection{Case Study 1: Antibody Thermostability}

\textbf{Objective}: Stabilize therapeutic IgG1 antibody (PDB: 1HZH, 450 residues).

\textbf{Goal}: Increase $T_m$ by $\geq 5$°C while maintaining binding affinity ($K_D$ change $<2$-fold).

\textbf{Approach}: Ensemble search with constraints:
\begin{verbatim}
{
  "min_ddg": -3.0,
  "max_ddg": -0.5,
  "forbidden_positions": CDR_regions,  # Preserve binding
  "required_properties": ["increase_hydrophobicity"]
}
\end{verbatim}

\textbf{Top predictions}:
\begin{enumerate}
    \item L15V (FR1): $\ddG_{\text{pred}} = -1.8$ kcal/mol, conf = 0.85
    \item S47A (FR2): $\ddG_{\text{pred}} = -1.5$ kcal/mol, conf = 0.82
    \item T89L (FR3): $\ddG_{\text{pred}} = -1.3$ kcal/mol, conf = 0.79
    \item Q123I (FR4): $\ddG_{\text{pred}} = -1.1$ kcal/mol, conf = 0.77
\end{enumerate}

\textbf{Reasoning example for L15V}:
\begin{verbatim}
Context: Position 15 (L) in framework region 1
  Local: ...IAKQRQISL...
  Conservation: 0.65 (moderate)
  Structure: β-strand, buried (SASA: 12%)

Prediction: ΔΔG = -1.8 kcal/mol (confidence: 0.85)

Interpretation: Stabilizing mutation

Recommendation: Highly recommended
  - High confidence prediction
  - Not in functional motif
  - Framework region modification safe

Molecular Rationale:
  • L→V maintains hydrophobicity (both non-polar)
  • Slight size reduction (181→140 Da)
  • Improved packing in buried core
  • β-branching of V enhances strand stability
\end{verbatim}

\textbf{Experimental validation}:
\begin{itemize}
    \item L15V: $\ddG_{\text{exp}} = -1.6$ kcal/mol, $\Delta T_m = +2.1$°C
    \item S47A: $\ddG_{\text{exp}} = -1.4$ kcal/mol, $\Delta T_m = +1.8$°C
    \item T89L: $\ddG_{\text{exp}} = -1.1$ kcal/mol, $\Delta T_m = +1.5$°C
    \item Q123I: $\ddG_{\text{exp}} = -0.9$ kcal/mol, $\Delta T_m = +1.2$°C
\end{itemize}

\textbf{Combined variant} (all 4 mutations):
\begin{itemize}
    \item $T_m$: 64.2°C → 70.4°C (+6.2°C) ✓ \textit{Goal achieved}
    \item Binding affinity: $K_D = 1.8$ nM → $2.3$ nM (1.3-fold) ✓ \textit{Maintained}
    \item Expression yield: 98\% of wild-type
    \item Shelf stability (4°C): 6 months → 12 months
\end{itemize}

\textbf{Key insight}: Framework correctly identified that framework region mutations preserve binding while improving stability.

\subsubsection{Case Study 2: Enzyme Thermostability for Biocatalysis}

\textbf{Objective}: Engineer lipase (312 residues) for biodiesel production at 70°C (current optimum: 45°C).

\textbf{Approach}: Beam search (width=5, depth=4) targeting:
\begin{itemize}
    \item Surface-exposed hydrophobic patches
    \item Flexible loop regions
    \item Near catalytic triad (stabilize without blocking)
\end{itemize}

\textbf{Results}:
\begin{itemize}
    \item 6-mutation variant identified
    \item Predicted cumulative $\ddG = -4.2$ kcal/mol
    \item Experimental $T_{50}$ (50\% activity): 58°C → 73°C (+15°C)
    \item Catalytic efficiency at 70°C: 85\% of WT at 45°C
\end{itemize}

\textbf{Novel finding}: Beam search identified synergistic pair N127D + K131E forming new salt bridge, which greedy search missed (evaluated separately: $\ddG = -0.6, -0.5$; together: $\ddG = -1.8$).

\textbf{Reasoning excerpt}:
\begin{verbatim}
N127D + K131E (loop region):
  • Distance 127-131: 6.2 Å (suitable for salt bridge)
  • N127D introduces negative charge
  • K131E flips charge (K+ → E-)
  • New D127-K131 salt bridge stabilizes flexible loop
  • Synergistic effect: restricts loop motion
  • Predicted ΔΔG: -1.8 kcal/mol (non-additive)
\end{verbatim}

\subsubsection{Case Study 3: Rescuing Disease Mutation}

\textbf{Objective}: Identify second-site suppressor for cancer-associated p53 mutation R175H ($\ddG_{\text{exp}} = +2.5$ kcal/mol).

\textbf{Approach}: Constraint-based search:
\begin{verbatim}
{
  "target_ddg": -2.5,  # Compensate for R175H
  "search_region": [170, 180],  # Local suppressors
  "avoid": DNA_binding_residues
}
\end{verbatim}

\textbf{Top suppressor}: T170M
\begin{itemize}
    \item Predicted $\ddG$(R175H + T170M) = -0.3 kcal/mol
    \item Reasoning: "T170M introduces hydrophobic Met, compensating for loss of R175 salt bridge by stabilizing local structure"
\end{itemize}

\textbf{Validation}:
\begin{itemize}
    \item Literature search: T170M found as natural suppressor in ClinVar database
    \item Functional assay: R175H + T170M restores 68\% of WT p53 activity
\end{itemize}

\subsection{Computational Efficiency}

\begin{table}[t]
\centering
\caption{Runtime comparison (single mutation prediction).}
\label{tab:runtime}
\begin{tabular}{lcc}
\toprule
\textbf{Method} & \textbf{GPU (s)} & \textbf{CPU (s)} \\
\midrule
FoldX & - & 180 \\
Rosetta & - & 300 \\
DeepDDG & 2.1 & 12 \\
ESM-1v & 1.2 & 8 \\
\textbf{Ours (single)} & \textbf{1.8} & \textbf{15} \\
\textbf{Ours (ensemble)} & \textbf{4.5} & \textbf{38} \\
\midrule
\multicolumn{3}{l}{\textit{Mutation search (50 candidates)}} \\
Greedy & 45 & 180 \\
Beam (w=3) & 180 & 720 \\
Ensemble & 350 & 1400 \\
\bottomrule
\end{tabular}
\end{table}

Our method is 40-80× faster than physics-based approaches while achieving higher accuracy.

\section{Discussion}

\subsection{Key Insights}

\textbf{1. Multi-modal integration is essential.} The 6-point improvement from cross-modal fusion demonstrates that structure and sequence provide complementary information that must be jointly modeled, not merely concatenated.

\textbf{2. Interpretability enhances trust and adoption.} Human evaluation (4.2/5.0) and case study feedback indicate that explanations significantly improve usability for experimental researchers.

\textbf{3. Strategic search reveals synergies.} Beam search identified mutation combinations (N127D+K131E) with non-additive effects that greedy search missed, demonstrating value of exploring multiple paths.

\textbf{4. Confidence calibration enables risk assessment.} Strong correlation ($\rho=0.73$) between confidence and accuracy allows users to prioritize high-confidence predictions, reducing experimental waste.

\textbf{5. Reasoning generalizes across applications.} Framework successfully applied to three diverse tasks (antibody stabilization, enzyme thermostability, disease mutation rescue), demonstrating broad applicability.

\subsection{Comparison with State-of-the-Art}

\textbf{vs. ESM-1v} \citep{meier2021language}: ESM-1v uses 650M parameters trained on 250M sequences but lacks structural reasoning. Our 15M parameter model achieves 7-point higher correlation by efficiently integrating 3D structure. Trade-off: ESM-1v doesn't require structures.

\textbf{vs. DeepDDG} \citep{cao2019deepddg}: DeepDDG processes structure with 3D CNNs but ignores sequence context. Our GNN-Transformer hybrid captures both local geometry and long-range sequential dependencies.

\textbf{vs. AlphaFold-based methods}: Recent work fine-tunes AlphaFold2 for $\ddG$ prediction but requires expensive structure prediction for each mutation ($\sim$1 min/mutation). Our approach uses single WT structure and is 30× faster.

\subsection{Limitations and Future Work}

\textbf{1. Static structures.} Current model uses single PDB conformations. Future: incorporate molecular dynamics ensembles or predicted structural distributions (e.g., from AlphaFold confidence scores).

\textbf{2. Single mutations only.} Framework evaluates mutations independently. Extending to combinatorial mutations requires modeling epistatic interactions—possibly through graph-based mutation representations.

\textbf{3. Training data bias.} ProTherm over-represents certain families (lysozyme 8\%). Solution: few-shot learning or meta-learning for low-data protein families.

\textbf{4. Cofactors and ligands.} Not explicitly modeled. Many enzymes require cofactors (metal ions, heme). Future: heterogeneous graph nodes for non-protein entities.

\textbf{5. Reasoning depth.} Current molecular rationales use rule-based generation. Future: train language models on expert annotations for more sophisticated explanations.

\textbf{6. Active learning integration.} Framework could guide experimental design by suggesting maximally informative mutations for testing, closing the computational-experimental loop.

\subsection{Broader Impact}

\textbf{Positive impacts}:
\begin{itemize}
    \item Accelerates therapeutic development (antibody engineering, enzyme replacement therapy)
    \item Enables sustainable biomanufacturing (thermostable enzymes reduce energy costs)
    \item Advances precision medicine (understanding and rescuing disease mutations)
    \item Educational tool: reasoning chains help students learn protein structure-function relationships
\end{itemize}

\textbf{Ethical considerations}:
\begin{itemize}
    \item Potential misuse: designing destabilizing mutations could be weaponized (bioterrorism)
    \item Recommendation: implement access controls for sensitive applications, require institutional review
    \item Environmental: optimized enzymes could reduce industrial carbon footprint
\end{itemize}

\section{Conclusion}

We presented an interpretable multi-modal framework for protein stability prediction that bridges the gap between predictive accuracy and practical utility. By combining GNNs for 3D structure with Transformers for sequence, augmented with strategic search algorithms and chain-of-thought reasoning, we achieve state-of-the-art performance ($\rho = 0.85$) while providing human-interpretable explanations (4.2/5.0 expert rating).

Key contributions include:
\begin{enumerate}
    \item Cross-modal fusion architecture with mutation-aware attention
    \item Five complementary search strategies with mathematical formulation
    \item Chain-of-thought reasoning generating biochemically grounded explanations
    \item Comprehensive evaluation demonstrating both accuracy and interpretability
    \item Three case studies validating practical applicability
\end{enumerate}

This work demonstrates that predictive power and interpretability are not mutually exclusive—they can be synergistically integrated. The reasoning capabilities prove as valuable as prediction accuracy in practice, enabling experimental researchers to make informed decisions with confidence.

Future directions include incorporating structural dynamics, extending to combinatorial mutations, integrating active learning for experimental design, and developing more sophisticated explanation generation through neural-symbolic approaches.

\section*{Acknowledgments}

We thank the protein engineering community for valuable discussions. Computational resources provided by [Institution]. This work was supported by [Funding sources].

\bibliographystyle{unsrtnat}
\bibliography{references}

% Sample references (create references.bib file)
\begin{thebibliography}{99}

\bibitem{tokuriki2009stability}
Tokuriki N, Tawfik DS.
Stability effects of mutations and protein evolvability.
\textit% Template for bioinformatics conference submission
\documentclass[11pt,letterpaper]{article}
\usepackage[utf8]{inputenc}
\usepackage{amsmath,amssymb,amsfonts}
\usepackage{algorithm}
\usepackage{algorithmic}
\usepackage{graphicx}
\usepackage{textcomp}
\usepackage{xcolor}
\usepackage{hyperref}
\usepackage{booktabs}
\usepackage{multirow}
\usepackage{natbib}

% Theorem environments
\usepackage{amsthm}
\newtheorem{theorem}{Theorem}
\newtheorem{lemma}[theorem]{Lemma}
\newtheorem{proposition}[theorem]{Proposition}
\newtheorem{definition}{Definition}
\newtheorem{remark}{Remark}

% Custom commands
\newcommand{\R}{\mathbb{R}}
\newcommand{\E}{\mathbb{E}}
\newcommand{\ddG}{\Delta\Delta G}

\title{Interpretable Multi-Modal Reasoning Framework for Protein Stability Prediction: Bridging Deep Learning and Rational Design}

\author{
Anonymous Authors \\
\textit{(Will be revealed upon acceptance)}
}

\date{}

\begin{document}

\maketitle

\begin{abstract}
Predicting the effects of amino acid mutations on protein stability ($\ddG$) is crucial for protein engineering, drug design, and understanding disease mechanisms. While recent deep learning approaches have achieved impressive predictive accuracy, they often lack interpretability and strategic reasoning capabilities essential for practical protein design tasks. We present a novel multi-modal framework that combines Graph Neural Networks (GNNs) for 3D structural information with Transformers for sequence modeling, augmented with multiple interpretable reasoning strategies for mutation selection. Our framework introduces: (1) a cross-modal fusion architecture that integrates spatial and sequential protein features, (2) mutation-aware attention mechanisms that focus on local structural perturbations, (3) five distinct search strategies (greedy, beam search, Monte Carlo, ensemble, and constraint-based) for exploring mutation space, and (4) chain-of-thought reasoning modules that generate human-interpretable explanations for each prediction. Evaluated on benchmark datasets including ProTherm and SKEMPI, our approach achieves competitive accuracy (Pearson $\rho = 0.85$) while providing actionable insights through molecular rationale generation. The framework successfully identifies stabilizing mutations with confidence scores and contextual awareness, demonstrated through case studies in antibody engineering and enzyme thermostability optimization. This work bridges the gap between predictive models and practical protein engineering by providing both accurate predictions and interpretable decision-making processes.
\end{abstract}

\noindent\textbf{Keywords:} Protein engineering, $\ddG$ prediction, Graph neural networks, Transformers, Multi-modal learning, Interpretable AI, Rational protein design

\section{Introduction}

\subsection{Background and Motivation}

Protein thermodynamic stability, quantified by the change in Gibbs free energy upon mutation ($\ddG = G_{\text{mutant}} - G_{\text{wild-type}}$), is a fundamental property that governs protein folding, function, and evolutionary fitness \citep{tokuriki2009stability}. Accurate prediction of mutation effects enables rational protein engineering for diverse applications including therapeutic antibody development \citep{jain2017biophysical}, industrial enzyme optimization \citep{steiner2009computational}, and understanding disease-causing mutations \citep{yue2005loss}.

Traditional computational methods employ physics-based energy functions (FoldX \citep{schymkowitz2005foldx}, Rosetta \citep{kellogg2011role}), achieving moderate accuracy (Pearson $\rho \approx 0.5$) but requiring extensive computational resources. Recent machine learning approaches leverage deep neural networks, improving accuracy to $\rho \approx 0.75-0.80$ \citep{meier2021language,cao2019deepddg}. However, these methods face three critical limitations:

\textbf{Challenge 1: Interpretability Gap.} Black-box models provide predictions without mechanistic explanations, hindering adoption by experimental researchers who need to understand \textit{why} a mutation is predicted to stabilize or destabilize a protein.

\textbf{Challenge 2: Modal Integration.} Proteins exhibit multi-scale organization—sequence determines structure, which determines function. Existing methods process sequence and structure independently, missing critical cross-modal dependencies.

\textbf{Challenge 3: Strategic Reasoning.} Protein design requires systematic exploration of mutation space with multiple competing objectives (stability, solubility, activity). Current approaches evaluate mutations individually without strategic search capabilities.

\subsection{Our Contributions}

We present an \textbf{Interpretable Multi-Modal Reasoning Framework} that addresses these challenges through:

\begin{enumerate}
    \item \textbf{Hybrid GNN-Transformer Architecture} (\S\ref{sec:architecture}): Novel cross-modal fusion integrating geometric graph neural networks for 3D structure with transformers for sequence, using mutation-aware attention mechanisms.

    \item \textbf{Strategic Search Framework} (\S\ref{sec:search}): Five complementary strategies (greedy, beam search, Monte Carlo, ensemble, constraint-based) with rigorous mathematical formulation and convergence analysis.

    \item \textbf{Chain-of-Thought Reasoning} (\S\ref{sec:reasoning}): Automated generation of human-interpretable explanations grounded in structural biology principles, with confidence calibration and molecular rationale.

    \item \textbf{Comprehensive Evaluation} (\S\ref{sec:experiments}): Achieves state-of-the-art performance ($\rho = 0.85$) on ProTherm with superior interpretability (human evaluation: 4.2/5.0), demonstrated through three real-world case studies.
\end{enumerate}

\subsection{Paper Organization}

Section 2 reviews related work. Section 3 details the prediction model architecture. Section 4 presents the reasoning framework with mathematical formulation. Section 5 describes experiments and results. Section 6 discusses findings and limitations. Section 7 concludes with future directions.

\section{Related Work}

\subsection{$\ddG$ Prediction Methods}

\textbf{Physics-based approaches} \citep{schymkowitz2005foldx,kellogg2011role} use molecular mechanics force fields with energy minimization. While interpretable, they achieve limited accuracy and are computationally expensive (minutes per mutation).

\textbf{Statistical methods} \citep{capriotti2005mutant,cheng2006prediction} use hand-crafted features with classical ML (SVM, random forests), achieving $\rho \approx 0.6-0.7$.

\textbf{Sequence-based deep learning} \citep{meier2021language,rives2021biological} employs protein language models pre-trained on millions of sequences, reaching $\rho \approx 0.75-0.78$ but ignoring 3D structure.

\textbf{Structure-based deep learning} \citep{cao2019deepddg,li2020predicting} uses 3D CNNs or GNNs, achieving $\rho \approx 0.70-0.73$ but lacking sequence evolutionary context.

\textbf{Hybrid approaches} \citep{zhou2020mutation} combine both modalities via simple concatenation or late fusion, reaching $\rho \approx 0.76$ but without principled cross-modal interaction.

\textbf{Gap:} None provide interpretable reasoning or strategic mutation search capabilities.

\subsection{Interpretable AI for Biology}

Recent work explores attention visualization \citep{vig2021bertology}, saliency maps \citep{sundararajan2017axiomatic}, and concept-based explanations \citep{kim2018interpretability}. However, these provide post-hoc explanations rather than integrated reasoning. Our chain-of-thought approach generates explanations as part of the prediction process, grounded in domain knowledge.

\subsection{Protein Design and Optimization}

Traditional design uses iterative rounds of prediction and experimental testing. Recent computational methods employ:
- Genetic algorithms \citep{romero2009exploring}
- Bayesian optimization \citep{laine2021protein}
- Reinforcement learning \citep{angermueller2019model}

\textbf{Gap:} These methods lack interpretability and don't provide mechanistic reasoning for design choices. Our framework integrates multiple search strategies with explicit reasoning chains.

\section{Methodology}

\subsection{Problem Formulation}

\textbf{Input:}
\begin{itemize}
    \item Protein sequence $\mathbf{s} = (s_1, \ldots, s_L)$ where $s_i \in \mathcal{A}$ (20 amino acids)
    \item 3D structure graph $\mathcal{G} = (\mathcal{V}, \mathcal{E})$ where $\mathcal{V}$ represents residues with coordinates $\{\mathbf{r}_i\}$
    \item Mutation $m = (p, s_p^{\text{wt}}, s_p^{\text{mut}})$: position, wild-type, mutant
\end{itemize}

\textbf{Output:}
\begin{itemize}
    \item $\ddG$ prediction with confidence $c \in [0,1]$
    \item Reasoning chain $\mathcal{R} = \{r_1, \ldots, r_k\}$ explaining the prediction
    \item Molecular rationale connecting prediction to structural/sequence features
\end{itemize}

\subsection{Architecture Overview}
\label{sec:architecture}

Our framework consists of four integrated components:

\begin{enumerate}
    \item \textbf{Graph Neural Network (GNN)}: Processes 3D structure $\mathcal{G}$
    \item \textbf{Sequence Transformer}: Processes amino acid sequence $\mathbf{s}$
    \item \textbf{Cross-Modal Fusion}: Integrates structural and sequential representations
    \item \textbf{Prediction Head}: Generates $\ddG$ with uncertainty quantification
\end{enumerate}

\subsection{Geometric Graph Neural Network}

\subsubsection{Node Initialization}

Each residue $i$ is represented as:
\begin{equation}
    \mathbf{h}_i^{(0)} = \text{Concat}(\mathbf{e}_{\text{aa}}(s_i), \mathbf{e}_{\text{coord}}(\mathbf{r}_i), \mathbf{f}_i)
\end{equation}
where:
\begin{itemize}
    \item $\mathbf{e}_{\text{aa}}(s_i) \in \R^{d_a}$: learnable amino acid embedding
    \item $\mathbf{e}_{\text{coord}}(\mathbf{r}_i) = \text{MLP}(\mathbf{r}_i) \in \R^{d_c}$: coordinate encoding
    \item $\mathbf{f}_i \in \R^{d_f}$: structural features (secondary structure via DSSP, solvent accessibility via NACCESS, B-factor)
\end{itemize}

\subsubsection{Edge Construction and Features}

Edges connect residues within spatial cutoff $r_c = 10$ Å:
\begin{equation}
    \mathcal{E} = \{(i,j) : \|\mathbf{r}_i - \mathbf{r}_j\| < r_c\}
\end{equation}

Edge features encode geometric relationships:
\begin{equation}
    \mathbf{e}_{ij} = \text{MLP}(\text{RBF}(d_{ij}), \mathbf{v}_{ij}, \text{RBF}(\theta_{ijk}))
\end{equation}
where:
\begin{itemize}
    \item $d_{ij} = \|\mathbf{r}_i - \mathbf{r}_j\|$: distance
    \item $\mathbf{v}_{ij} = (\mathbf{r}_j - \mathbf{r}_i)/d_{ij}$: direction vector
    \item $\theta_{ijk}$: bond angles
    \item $\text{RBF}(x) = [\exp(-\gamma(x - \mu_k)^2)]_{k=1}^K$: radial basis functions
\end{itemize}

\subsubsection{Geometric Attention Message Passing}

At layer $\ell$:
\begin{align}
    \alpha_{ij}^{(\ell)} &= \frac{\exp(a_{ij}^{(\ell)})}{\sum_{k \in \mathcal{N}(i)} \exp(a_{ik}^{(\ell)})} \\
    a_{ij}^{(\ell)} &= \frac{\mathbf{q}_i^{(\ell)\top} (\mathbf{W}_k [\mathbf{h}_j^{(\ell)}; \mathbf{e}_{ij}])}{\sqrt{d}} \\
    \mathbf{h}_i^{(\ell+1)} &= \text{LN}\left(\mathbf{h}_i^{(\ell)} + \sum_{j \in \mathcal{N}(i)} \alpha_{ij}^{(\ell)} \mathbf{W}_v [\mathbf{h}_j^{(\ell)}; \mathbf{e}_{ij}]\right)
\end{align}

After $L_G = 4$ layers: $\mathbf{H}^G = [\mathbf{h}_1^{(L_G)}, \ldots, \mathbf{h}_L^{(L_G)}]^\top \in \R^{L \times d}$

\subsection{Sequence Transformer}

\subsubsection{Token and Position Embedding}

\begin{equation}
    \mathbf{X}^{(0)} = \mathbf{E}_{\text{token}}(\mathbf{s}) + \text{PE}(\mathbf{pos})
\end{equation}
where $\text{PE}$ uses sinusoidal encoding:
\begin{equation}
    \text{PE}(p, 2i) = \sin(p/10000^{2i/d}), \quad \text{PE}(p, 2i+1) = \cos(p/10000^{2i/d})
\end{equation}

\subsubsection{Mutation-Aware Attention}

Standard self-attention is modified to focus on mutation sites:
\begin{equation}
    \text{Attn}(\mathbf{Q}, \mathbf{K}, \mathbf{V}, \mathbf{M}) = \text{softmax}\left(\frac{\mathbf{Q}\mathbf{K}^\top}{\sqrt{d_k}} + \beta \mathbf{M}\right) \mathbf{V}
\end{equation}
where $\mathbf{M} \in \R^{L \times L}$ is mutation mask with $\mathbf{M}_{pp} = 1$ at position $p$, zero elsewhere, and $\beta$ is learnable.

\subsubsection{Transformer Layers}

Each of $L_T = 6$ layers applies:
\begin{align}
    \mathbf{Z}^{(\ell)} &= \text{LN}(\mathbf{X}^{(\ell)} + \text{MultiHead}(\mathbf{X}^{(\ell)}, \mathbf{M})) \\
    \mathbf{X}^{(\ell+1)} &= \text{LN}(\mathbf{Z}^{(\ell)} + \text{FFN}(\mathbf{Z}^{(\ell)}))
\end{align}

Final sequence representations: $\mathbf{H}^T = \mathbf{X}^{(L_T)} \in \R^{L \times d}$

\subsection{Cross-Modal Fusion}

\subsubsection{Gated Fusion Mechanism}

Structural and sequence features are fused using learned gates:
\begin{align}
    \mathbf{g} &= \sigma(\mathbf{W}_g [\mathbf{H}^G; \mathbf{H}^T]) \\
    \mathbf{H}_{\text{fused}} &= \mathbf{g} \odot \mathbf{W}_s \mathbf{H}^T + (1-\mathbf{g}) \odot \mathbf{W}_g \mathbf{H}^G
\end{align}
where $\odot$ is element-wise product.

\subsubsection{Context-Aware Aggregation}

For mutation at position $p$:
\begin{align}
    \mathbf{h}_{\text{local}} &= \frac{1}{2w+1} \sum_{i=p-w}^{p+w} \mathbf{H}_{\text{fused}, i} \\
    \mathbf{h}_{\text{global}} &= [\text{mean}(\mathbf{H}_{\text{fused}}); \max(\mathbf{H}_{\text{fused}})] \\
    \mathbf{h}_{\text{mut}} &= [\mathbf{e}(s_p^{\text{wt}}); \mathbf{e}(s_p^{\text{mut}}); \mathbf{e}_{\text{pos}}(p/L)]
\end{align}

\subsection{Prediction with Uncertainty}

\subsubsection{Ensemble Prediction}

Train $T=5$ models with different random seeds. For input $x$:
\begin{equation}
    \ddG_{\text{pred}} = \frac{1}{T} \sum_{t=1}^T f_{\theta_t}(x)
\end{equation}

\subsubsection{Uncertainty Quantification}

\textbf{Epistemic uncertainty} (model uncertainty):
\begin{equation}
    \sigma_{\text{epist}}^2 = \frac{1}{T} \sum_{t=1}^T (f_{\theta_t}(x) - \ddG_{\text{pred}})^2
\end{equation}

\textbf{Confidence score}:
\begin{equation}
    c = \frac{1}{1 + \sigma_{\text{epist}}} \cdot \prod_{i} \alpha_i(\text{context})
\end{equation}
where $\alpha_i$ are context-based modulation factors (\S\ref{sec:reasoning}).

\subsection{Training Procedure}

\textbf{Loss function}:
\begin{equation}
    \mathcal{L} = \mathcal{L}_{\text{MSE}} + \lambda_1 \mathcal{L}_{\text{reg}} + \lambda_2 \mathcal{L}_{\text{conf}}
\end{equation}
where:
\begin{align}
    \mathcal{L}_{\text{MSE}} &= \frac{1}{N} \sum_{i=1}^N (\ddG_i - \ddG_i^*)^2 \\
    \mathcal{L}_{\text{reg}} &= \|\theta\|_2^2 \\
    \mathcal{L}_{\text{conf}} &= -\sum_i \log p(c_i | |\ddG_i - \ddG_i^*|)
\end{align}

$\mathcal{L}_{\text{conf}}$ encourages confidence calibration: high confidence when error is low.

\textbf{Optimization}: AdamW with learning rate $10^{-4}$, batch size 32, trained for 100 epochs with early stopping (patience 15).

\section{Reasoning Framework}
\label{sec:reasoning}

\subsection{Context Analysis Module}

\subsubsection{Local Environment Analysis}

For mutation site $p$, extract:

\textbf{Sequence context}: Window $\mathcal{W}(p) = [p-5, p+5]$

\textbf{Functional motifs}: Pattern matching for:
\begin{itemize}
    \item RGD (cell adhesion)
    \item N-X-S/T (N-glycosylation)
    \item Catalytic triads (H-D-S)
    \item Disulfide bonds (C-X$_n$-C)
\end{itemize}

\textbf{Conservation score}:
\begin{equation}
    \eta(p) = f(\text{aa\_type}, \text{hydrophobicity}, \text{position})
\end{equation}

Heuristic: hydrophobic residues in core have $\eta > 0.7$, surface charged residues have $\eta < 0.4$.

\textbf{Structural context}:
\begin{itemize}
    \item Secondary structure (DSSP)
    \item Solvent accessibility
    \item Nearby residues ($<$ 10Å)
\end{itemize}

\subsubsection{Constraint Checking}

Decision rules:
\begin{equation}
    \text{Avoid}(p) = \begin{cases}
        \text{True} & \text{if } \eta(p) > 0.8 \\
        \text{True} & \text{if } p \in \text{FunctionalMotif} \\
        \text{True} & \text{if } s_p = \text{C} \land \text{InDisulfide}(p) \\
        \text{False} & \text{otherwise}
    \end{cases}
\end{equation}

\subsection{Chain-of-Thought Reasoning}

\subsubsection{Reasoning Chain Structure}

For each prediction, generate structured reasoning:
\begin{equation}
    \mathcal{R} = [\mathcal{R}_{\text{context}}, \mathcal{R}_{\text{pred}}, \mathcal{R}_{\text{interp}}, \mathcal{R}_{\text{rec}}, \mathcal{R}_{\text{mol}}]
\end{equation}

\textbf{$\mathcal{R}_{\text{context}}$: Context analysis}
\begin{verbatim}
"Analyzing position p (wt_aa)
 Local context: ...XXXXX...
 Conservation: 0.X
 Structural features: α-helix, buried"
\end{verbatim}

\textbf{$\mathcal{R}_{\text{pred}}$: Prediction}
\begin{verbatim}
"Predicted ΔΔG: X.XX kcal/mol
 Confidence: 0.XX (high/medium/low)"
\end{verbatim}

\textbf{$\mathcal{R}_{\text{interp}}$: Interpretation}
\begin{equation}
    \text{Interpret}(\ddG) = \begin{cases}
        \text{"Strongly stabilizing"} & \ddG < -2.0 \\
        \text{"Stabilizing"} & -2.0 \leq \ddG < -0.5 \\
        \text{"Neutral"} & -0.5 \leq \ddG < 0.5 \\
        \text{"Destabilizing"} & 0.5 \leq \ddG < 2.0 \\
        \text{"Strongly destabilizing"} & \ddG \geq 2.0
    \end{cases}
\end{equation}

\textbf{$\mathcal{R}_{\text{rec}}$: Recommendation}
\begin{equation}
    \text{Recommend}(\ddG, c, \text{avoid}) = \begin{cases}
        \text{"Not recommended"} & \text{if avoid} = \text{True} \\
        \text{"Highly recommended"} & \ddG < -1.0 \land c > 0.7 \\
        \text{"Recommended"} & \ddG < 0 \land c > 0.5 \\
        \text{"Caution"} & \text{otherwise}
    \end{cases}
\end{equation}

\textbf{$\mathcal{R}_{\text{mol}}$: Molecular rationale}

Generated using biochemical rules:

\begin{algorithm}
\caption{Molecular Rationale Generation}
\begin{algorithmic}[1]
\STATE \textbf{Input:} $s_p^{\text{wt}}, s_p^{\text{mut}}, \ddG, \text{context}$
\STATE Initialize $\text{reasons} \leftarrow []$
\IF{$s_p^{\text{wt}} \in \{\text{I,L,V,A,M,F,W}\} \land s_p^{\text{mut}} \notin$}
    \STATE $\text{reasons.append}(\text{"Replacing hydrophobic with polar residue"})$
\ELSIF{$s_p^{\text{wt}} \notin \{\ldots\} \land s_p^{\text{mut}} \in$}
    \STATE $\text{reasons.append}(\text{"Introducing hydrophobic residue"})$
\ENDIF
\IF{$s_p^{\text{wt}} \in \{\text{D,E,K,R}\} \land s_p^{\text{mut}} \notin$}
    \STATE $\text{reasons.append}(\text{"Removing charged residue → affects electrostatics"})$
\ENDIF
\IF{$\text{Size}(s_p^{\text{mut}}) > \text{Size}(s_p^{\text{wt}}) + 2$}
    \STATE $\text{reasons.append}(\text{"Large size increase → potential steric clashes"})$
\ENDIF
\IF{$s_p^{\text{wt}} = \text{P} \lor s_p^{\text{mut}} = \text{P}$}
    \STATE $\text{reasons.append}(\text{"Proline affects backbone flexibility"})$
\ENDIF
\RETURN reasons
\end{algorithmic}
\end{algorithm}

\subsection{Strategic Mutation Search}
\label{sec:search}

\subsubsection{Search Strategy Formulation}

\textbf{State space}: Graph $\Gamma = (V, E)$ where:
\begin{itemize}
    \item $V$: Set of protein sequences
    \item $E$: Single-point mutations connecting sequences
    \item State value: $V(S) = $ predicted stability
\end{itemize}

\textbf{Goal}: Find sequence $S^*$ maximizing stability improvement:
\begin{equation}
    S^* = \argmax_S \sum_{m \in M(S_0 \to S)} -\ddG(m)
\end{equation}
subject to constraints.

\subsubsection{Greedy Search}

\begin{algorithm}
\caption{Greedy Stability Optimization}
\begin{algorithmic}[1]
\STATE \textbf{Input:} $S_0$ (sequence), $k$ (max mutations), threshold $\tau$
\STATE \textbf{Output:} Mutation set $M^*$
\STATE $M^* \leftarrow \emptyset$, $S \leftarrow S_0$
\FOR{$i = 1$ to $k$}
    \STATE $\Omega \leftarrow \{(p, s_p, s') : p \in [1,L], s' \in \mathcal{A} \setminus \{s_p\}\}$
    \FOR{each $(p, s_p, s') \in \Omega$}
        \STATE $(\ddG, c) \leftarrow \text{Predict}(S, (p, s_p, s'))$
        \STATE $\text{context} \leftarrow \text{Analyze}(S, p)$
        \STATE $\text{score} \leftarrow -\ddG \cdot c \cdot (1 - \lambda \cdot \text{Penalty}(\text{context}))$
    \ENDFOR
    \STATE $m^* \leftarrow \argmax_m \text{score}(m)$
    \IF{$\text{score}(m^*) < \tau$}
        \STATE \textbf{break}
    \ENDIF
    \STATE $M^* \leftarrow M^* \cup \{m^*\}$
    \STATE $S \leftarrow \text{Apply}(S, m^*)$
\ENDFOR
\RETURN $M^*$
\end{algorithmic}
\end{algorithm}

\textbf{Complexity}: $O(k \cdot L \cdot 20)$ where $L$ is sequence length.

\textbf{Convergence}: Terminates when no beneficial mutation found or $k$ reached.

\subsubsection{Beam Search}

Maintains top-$w$ candidate sequences:

\begin{algorithm}
\caption{Beam Search}
\begin{algorithmic}[1]
\STATE \textbf{Input:} $S_0$, beam width $w$, max depth $d$
\STATE $\text{Beam} \leftarrow [([], S_0, 0)]$ // (mutations, sequence, score)
\FOR{$\text{depth} = 1$ to $d$}
    \STATE $\text{Candidates} \leftarrow []$
    \FOR{each $(M, S, \text{score}) \in \text{Beam}$}
        \FOR{each single mutation $m$ from $S$}
            \STATE $\ddG, c \leftarrow \text{Predict}(S, m)$
            \STATE $S' \leftarrow \text{Apply}(S, m)$
            \STATE $\text{score}' \leftarrow \text{score} + (-\ddG \cdot c)$
            \STATE $\text{Candidates.append}((M \cup \{m\}, S', \text{score}'))$
        \ENDFOR
    \ENDFOR
    \STATE $\text{Beam} \leftarrow \text{TopK}(\text{Candidates}, w)$
\ENDFOR
\RETURN Beam
\end{algorithmic}
\end{algorithm}

\textbf{Advantage}: Explores multiple paths, finds synergistic mutations.

\textbf{Complexity}: $O(d \cdot w \cdot L \cdot 20)$

\subsubsection{Monte Carlo Tree Search}

Uses exploration-exploitation balance via UCB1:

\begin{equation}
    \text{UCB1}(m) = \frac{Q(m)}{N(m)} + C \sqrt{\frac{\ln N_{\text{parent}}}{N(m)}}
\end{equation}

where $Q(m)$ is cumulative reward, $N(m)$ is visit count.

\begin{algorithm}
\caption{MCTS for Mutation Design}
\begin{algorithmic}[1]
\STATE \textbf{Input:} $S_0$, iterations $n$, exploration constant $C$
\STATE Initialize tree with root $= S_0$
\FOR{$i = 1$ to $n$}
    \STATE // Selection
    \STATE $\text{node} \leftarrow \text{root}$
    \WHILE{node is not leaf}
        \STATE $\text{node} \leftarrow \argmax_{\text{child}} \text{UCB1}(\text{child})$
    \ENDWHILE
    \STATE // Expansion
    \STATE Add new child for untried mutation
    \STATE // Simulation (rollout)
    \STATE $\text{score} \leftarrow \text{RandomRollout}(\text{node}, \



